\documentclass[a4paper,12pt]{article}
\usepackage[utf8]{inputenc}
\usepackage[T1]{fontenc}
\usepackage{amsmath}
\usepackage{amssymb}
\usepackage{amsthm}
\usepackage{geometry}
\usepackage{graphicx}
\usepackage{xcolor}
\usepackage{listings}
\usepackage{algorithm}
\usepackage{algpseudocode}
\usepackage{float}
\floatname{algorithm}{Algorithmus}

% Deutsche Schlüsselwörter für algpseudocode
\algrenewcommand\algorithmicif{\textbf{wenn}}
\algrenewcommand\algorithmicelse{\textbf{sonst}}
\algrenewcommand\algorithmicthen{\textbf{dann}}
\algrenewcommand\algorithmicend{\textbf{Ende}}
\algrenewcommand\algorithmicfor{\textbf{für}}
\algrenewcommand\algorithmicforall{\textbf{für alle}}
\algrenewcommand\algorithmicdo{\textbf{tue}}
\algrenewcommand\algorithmicwhile{\textbf{solange}}
\algrenewcommand\algorithmicrepeat{\textbf{wiederhole}}
\algrenewcommand\algorithmicuntil{\textbf{bis}}
\algrenewcommand\algorithmicreturn{\textbf{gib zurück}}
\algrenewcommand\algorithmicfunction{\textbf{Prozedur}}
\algrenewcommand\algorithmicprocedure{\textbf{Prozedur}}
\algrenewcommand\algorithmicrequire{\textbf{Eingabe:}}
\algrenewcommand\algorithmicensure{\textbf{Ausgabe:}}
\algrenewcommand\algorithmiccomment[1]{\hfill\(\triangleright\) #1}

% Neue deutsche Block-Befehle als Aliase
\algnewcommand\Prozedur[2]{\Procedure{#1}{#2}}
\algnewcommand\EndProzedur{\EndProcedure}
\algnewcommand\Waehrend[1]{\While{#1}}
\algnewcommand\EndeWaehrend{\EndWhile}
\algnewcommand\Fuer[1]{\For{#1}}
\algnewcommand\FuerAlle[1]{\ForAll{#1}}
\algnewcommand\EndeFuer{\EndFor}
\algnewcommand\Wenn[1]{\If{#1}}
\algnewcommand\SonstWenn[1]{\ElsIf{#1}}
\algnewcommand\Sonst{\Else}
\algnewcommand\EndeWenn{\EndIf}
\algnewcommand\Wiederhole{\Repeat}
\algnewcommand\Bis[1]{\Until{#1}}
\algnewcommand\Eingabe{\Require}
\algnewcommand\Ausgabe{\Ensure}
\usepackage{chngcntr}
\usepackage{booktabs}

% Hyperlinkfarben aus subgraph.tex (RGB 25,70,140 = dunkles Blau)
\definecolor{mainblue}{RGB}{25, 70, 140}
\definecolor{accentred}{RGB}{180, 50, 30}

\usepackage{hyperref}
\hypersetup{
	colorlinks=true,
	linkcolor=mainblue,
	citecolor=accentred,
	urlcolor=mainblue
}

\geometry{margin=2.5cm}
\usepackage[protrusion=true,expansion=false]{microtype}       % Besserer Mikrotypographie-Ausgleich
\usepackage{rotating}        % Querformat-Tabellen mit sidewaystable

% Code-Highlighting für PDDL
\lstdefinelanguage{PDDL}{
	keywords={define, domain, problem, requirements, types, predicates, action, parameters, precondition, effect, and, not, or, forall, exists, when, either, objects, init, goal},
	keywordstyle=\color{blue}\bfseries,
	comment=[l]{;},
	commentstyle=\color{gray}\itshape,
	string=[b]",
	stringstyle=\color{red},
	morestring=[b]',
	sensitive=false
}

\lstset{
	basicstyle=\ttfamily\small,
	keywordstyle=\color{blue},
	commentstyle=\color{gray},
	stringstyle=\color{red},
	numbers=left,
	numberstyle=\tiny,
	breaklines=true,
	frame=single,
	showstringspaces=false
}

% Deutsche Bezeichnungen für Verzeichnisse und Gleitobjekte
\renewcommand{\contentsname}{Inhaltsverzeichnis}
\renewcommand{\listfigurename}{Abbildungsverzeichnis}
\renewcommand{\listtablename}{Tabellenverzeichnis}
\renewcommand{\figurename}{Abbildung}
\renewcommand{\tablename}{Tabelle}
\renewcommand{\refname}{Literaturverzeichnis}

\renewcommand{\abstractname}{Zusammenfassung}
\renewcommand{\proofname}{Beweis}

\theoremstyle{definition}
\newtheorem{definition}{Definition}[section]
\newtheorem{beispiel}{Beispiel}[section]
\newtheorem{theorem}{Theorem}[section]
\newtheorem{satz}{Satz}[section]
\newtheorem{lemma}{Lemma}[section]
\newtheorem{example}{Beispiel}[section]
\newtheorem{remark}{Bemerkung}[section]

\title{Temporales Model Checking und PDDL Planning\\
	\large{Formale Verifikation von Planungssystemen}}

\author{Stephan Epp}

\date{\today}

\begin{document}
	
\maketitle

\tableofcontents
\newpage

\section{Einleitung}

Temporales Model Checking ist eine fundamentale Technik zur Verifikation zeitabhängiger Eigenschaften von Systemen. Während klassisches Model Checking primär strukturelle Eigenschaften von Systemen untersucht, konzentriert sich das temporale Model Checking auf die zeitliche Entwicklung von Systemzuständen und die Gültigkeit temporaler Logikformeln über Ausführungspfaden.

Diese Arbeit führt in die Grundlagen des temporalen Model Checkings ein und demonstriert die praktische Anwendung anhand von PDDL (Planning Domain Definition Language) Plannern. PDDL ist ein standardisierter Formalismus zur Beschreibung von Planungsproblemen, der in der künstlichen Intelligenz und automatisierten Planung weit verbreitet ist.

Der Beitrag dieser Arbeit besteht in: (1.) Einführung in temporale Logiken LTL und CTL, (2.) Algorithmen für temporales Model Checking, (3.) Formalisierung von PDDL Planungsproblemen, (4.) Praktische Beispiele mit konkreten Domänen und Probleminstanzen, (5.) Verifikation von Planeigenschaften mittels temporaler Logik, (6.) Der Subgraph Algorithmus als polynomiale Alternative mit $O(n^3)$-Laufzeit sowie seine Integration in den gesamten Verifikations-Workflow.

\subsection{Motivation}

Planungssysteme werden in sicherheitskritischen Bereichen eingesetzt, darunter Robotersteuerung, Logistik und automatisierte Entscheidungssysteme. Die Korrektheit solcher Systeme ist essentiell. Temporales Model Checking bietet die Möglichkeit, formale Garantien über das Verhalten von Planern zu geben:

\begin{itemize}
	\item[-] \textbf{Safety}: Werden niemals unerwünschte Zustände erreicht?
	\item[-] \textbf{Liveness}: Werden erwünschte Ziele schließlich erreicht?
	\item[-] \textbf{Fairness}: Werden Ressourcen gerecht verteilt?
	\item[-] \textbf{Deadlock-Freiheit}: Kann das System stets fortschreiten?
\end{itemize}

Die Kombination von PDDL Planning mit temporalem Model Checking ermöglicht die automatisierte Verifikation von Plänen und Domänenbeschreibungen.

\section{Grundlagen Temporaler Logiken}

Temporale Logiken erweitern klassische Aussagenlogik um Operatoren zur Beschreibung zeitlicher Abläufe. Die beiden wichtigsten temporalen Logiken sind Linear Temporal Logic (LTL) und Computation Tree Logic (CTL).

\subsection{Linear Temporal Logic (LTL)}

LTL beschreibt Eigenschaften über lineare Ausführungspfade eines Systems. Ein Pfad ist eine unendliche Sequenz von Zuständen $s_0, s_1, s_2, \ldots$

\subsubsection{Syntax}

Die Syntax von LTL ist induktiv definiert:
\begin{align}
\varphi ::= \text{true} \mid p \mid \neg \varphi \mid \varphi_1 \land \varphi_2 \mid \mathbf{X}\varphi \mid \mathbf{F}\varphi \mid \mathbf{G}\varphi \mid \varphi_1 \mathbf{U} \varphi_2
\end{align}

Dabei sind:
\begin{itemize}
	\item[-] $p$ eine atomare Proposition (Zustandsvariable)
	\item[-] $\mathbf{X}\varphi$ (neXt): $\varphi$ gilt im nächsten Zustand
	\item[-] $\mathbf{F}\varphi$ (Finally): $\varphi$ gilt irgendwann in der Zukunft
	\item[-] $\mathbf{G}\varphi$ (Globally): $\varphi$ gilt immer
	\item[-] $\varphi_1 \mathbf{U} \varphi_2$ (Until): $\varphi_1$ gilt, bis $\varphi_2$ gilt
\end{itemize}

\begin{beispiel}[LTL Formeln]
	\begin{enumerate}
		\item $\mathbf{G}(\text{request} \to \mathbf{F}\text{grant})$: Jede Anfrage wird schließlich bearbeitet
		\item $\mathbf{G}\neg\text{deadlock}$: Es gibt niemals einen Deadlock
		\item $\mathbf{F}\mathbf{G}\text{stable}$: Das System stabilisiert sich schließlich
		\item $\mathbf{G}(\text{button\_pressed} \to \mathbf{X}\text{light\_on})$: Nach Knopfdruck geht sofort das Licht an
	\end{enumerate}
\end{beispiel}

\subsubsection{Semantik}

Die Semantik von LTL wird über Pfade $\pi = s_0, s_1, s_2, \ldots$ definiert, wobei $\pi^i$ den Suffix-Pfad ab Position $i$ bezeichnet:

\begin{align}
\pi \models p &\iff p \in L(s_0) \\
\pi \models \neg\varphi &\iff \pi \not\models \varphi \\
\pi \models \varphi_1 \land \varphi_2 &\iff \pi \models \varphi_1 \text{ und } \pi \models \varphi_2 \\
\pi \models \mathbf{X}\varphi &\iff \pi^1 \models \varphi \\
\pi \models \mathbf{F}\varphi &\iff \exists i \geq 0: \pi^i \models \varphi \\
\pi \models \mathbf{G}\varphi &\iff \forall i \geq 0: \pi^i \models \varphi \\
\pi \models \varphi_1 \mathbf{U} \varphi_2 &\iff \exists j \geq 0: \pi^j \models \varphi_2 \text{ und } \forall 0 \leq i < j: \pi^i \models \varphi_1
\end{align}

\subsection{Computation Tree Logic (CTL)}

CTL beschreibt Eigenschaften über Berechnungsbäume, die alle möglichen Ausführungen eines Systems repräsentieren. Im Gegensatz zu LTL quantifiziert CTL explizit über Pfade.

\subsubsection{Syntax}

Die Syntax von CTL kombiniert Pfadquantoren mit temporalen Operatoren:
\begin{multline}
\varphi ::= \text{true} \mid p \mid \neg \varphi \mid \varphi_1 \land \varphi_2
  \mid \mathbf{AX}\varphi \mid \mathbf{EX}\varphi \mid \mathbf{AF}\varphi \mid \mathbf{EF}\varphi \\
  \mid \mathbf{AG}\varphi \mid \mathbf{EG}\varphi
  \mid \mathbf{A}[\varphi_1 \mathbf{U} \varphi_2] \mid \mathbf{E}[\varphi_1 \mathbf{U} \varphi_2]
\end{multline}

Dabei sind:
\begin{itemize}
	\item[-] $\mathbf{A}$ (All paths): Auf allen Pfaden gilt..
	\item[-] $\mathbf{E}$ (Exists path): Es existiert ein Pfad, auf dem gilt..
	\item[-] $\mathbf{X}, \mathbf{F}, \mathbf{G}, \mathbf{U}$ wie in LTL
\end{itemize}

\begin{beispiel}[CTL Formeln]
	\begin{enumerate}
		\item $\mathbf{AG}(\text{request} \to \mathbf{AF}\text{grant})$: Auf allen Pfaden wird jede Anfrage schließlich bearbeitet
		\item $\mathbf{EF}\text{goal}$: Es existiert ein Pfad zum Zielzustand
		\item $\mathbf{AG}\neg\text{deadlock}$: Auf allen Pfaden gibt es niemals Deadlock
		\item $\mathbf{EG}\text{alive}$: Es existiert ein Pfad, auf dem das System ewig läuft
	\end{enumerate}
\end{beispiel}

\subsubsection{Semantik}

Die Semantik von CTL wird über einen Kripke-Strukturen $M = (S, R, L)$ definiert:

\begin{align}
M, s \models p &\iff p \in L(s) \\
M, s \models \neg\varphi &\iff M, s \not\models \varphi \\
M, s \models \varphi_1 \land \varphi_2 &\iff M, s \models \varphi_1 \text{ und } M, s \models \varphi_2 \\
M, s \models \mathbf{AX}\varphi &\iff \forall s': (s,s') \in R \to M, s' \models \varphi \\
M, s \models \mathbf{EX}\varphi &\iff \exists s': (s,s') \in R \land M, s' \models \varphi \\
M, s \models \mathbf{AG}\varphi &\iff \forall \text{ Pfade } \pi \text{ ab } s, \forall i: M, \pi_i \models \varphi \\
M, s \models \mathbf{EF}\varphi &\iff \exists \text{ Pfad } \pi \text{ ab } s, \exists i: M, \pi_i \models \varphi
\end{align}

\subsection{Vergleich LTL und CTL}

Tabelle \ref{tab:ltl_ctl_comparison} zeigt die wichtigsten Unterschiede:

\begin{table}[h]
	\centering
	\caption{Vergleich von LTL und CTL}
	\label{tab:ltl_ctl_comparison}
	\begin{tabular}{lll}
		\toprule
		Aspekt & LTL & CTL \\
		\midrule
		Quantifizierung & Implizit über alle Pfade & Explizit (A, E) \\
		Ausdrucksstärke & Linear & Verzweigend \\
		Typische Eigenschaften & Fairness, Reaktivität & Erreichbarkeit, Invarianten \\
		Komplexität MC & PSPACE-vollständig & P (für CTL) \\
		Anwendung & Spezifikation & Verifikation \\
		\bottomrule
	\end{tabular}
\end{table}

\begin{remark}
	LTL und CTL sind nicht gegenseitig ausdrucksstark. Es gibt LTL-Formeln, die nicht in CTL ausdrückbar sind (z.B. $\mathbf{FG}p$) und umgekehrt (z.B. $\mathbf{AG}\mathbf{EF}p$). CTL* vereint beide Logiken.
\end{remark}

\section{Algorithmen für Temporales Model Checking}

Dieser Abschnitt präsentiert die grundlegenden Algorithmen für LTL und CTL Model Checking.

\subsection{CTL Model Checking}

Der CTL Model Checking Algorithmus arbeitet durch rekursive Berechnung von Zustandsmengen, die eine Formel erfüllen.

\subsubsection{Der SAT-Algorithmus}

Gegeben eine Kripke-Struktur $M = (S, R, L)$ und eine CTL-Formel $\varphi$, berechnet der SAT-Algorithmus die Menge $\text{SAT}(\varphi) = \{s \in S \mid M, s \models \varphi\}$.

\begin{algorithm}[H]
	\caption{CTL-Modellprüfung (SAT-Algorithmus)}
	\begin{algorithmic}[1]
		\Prozedur{SAT}{$M, \varphi$}
		\Wenn{$\varphi = \text{true}$}
		\State \Return $S$
		\SonstWenn{$\varphi = p$}
		\State \Return $\{s \in S \mid p \in L(s)\}$
		\SonstWenn{$\varphi = \neg\psi$}
		\State \Return $S \setminus \text{SAT}(M, \psi)$
		\SonstWenn{$\varphi = \psi_1 \land \psi_2$}
		\State \Return $\text{SAT}(M, \psi_1) \cap \text{SAT}(M, \psi_2)$
		\SonstWenn{$\varphi = \mathbf{EX}\psi$}
		\State \Return $\{s \in S \mid \exists s': (s,s') \in R \land s' \in \text{SAT}(M, \psi)\}$
		\SonstWenn{$\varphi = \mathbf{EG}\psi$}
		\State \Return $\text{EG-SAT}(M, \psi)$
		\SonstWenn{$\varphi = \mathbf{E}[\psi_1 \mathbf{U} \psi_2]$}
		\State \Return $\text{EU-SAT}(M, \psi_1, \psi_2)$
		\EndeWenn
		\EndProzedur
	\end{algorithmic}
\end{algorithm}

Die Hilfsfunktionen EG-SAT und EU-SAT werden durch Fixpunkt-Berechnung implementiert.

\subsubsection{EG-SAT (Eventually Globally)}

\begin{algorithm}[H]
	\caption{EG-SAT-Berechnung}
	\begin{algorithmic}[1]
		\Prozedur{EG-SAT}{$M, \psi$}
		\State $T \gets \text{SAT}(M, \psi)$ \Comment{Starte mit allen $\psi$-Zuständen}
		\State $T' \gets \emptyset$
		\Waehrend{$T \neq T'$}
		\State $T' \gets T$
		\State $T \gets T \cap \{s \mid \exists s': (s,s') \in R \land s' \in T\}$ \Comment{Entferne Zustände ohne Nachfolger in $T$}
		\EndeWaehrend
		\State \Return $T$
		\EndProzedur
	\end{algorithmic}
\end{algorithm}

\subsubsection{EU-SAT (Exists Until)}

\begin{algorithm}[H]
	\caption{EU-SAT-Berechnung}
	\begin{algorithmic}[1]
		\Prozedur{EU-SAT}{$M, \psi_1, \psi_2$}
		\State $T \gets \text{SAT}(M, \psi_2)$ \Comment{Starte mit $\psi_2$-Zuständen}
		\State $T' \gets \emptyset$
		\Waehrend{$T \neq T'$}
		\State $T' \gets T$
		\State $T \gets T \cup (\text{SAT}(M, \psi_1) \cap \{s \mid \exists s': (s,s') \in R \land s' \in T\})$
		\EndeWaehrend
		\State \Return $T$
		\EndProzedur
	\end{algorithmic}
\end{algorithm}

\begin{theorem}[Korrektheit CTL Model Checking]
	Der SAT-Algorithmus terminiert für alle CTL-Formeln $\varphi$ und berechnet korrekt die Menge $\text{SAT}(\varphi)$.
\end{theorem}

\begin{proof}
	Die Terminierung folgt aus der Monotonie der Fixpunkt-Operationen und der Endlichkeit des Zustandsraums $S$. Die Korrektheit ergibt sich durch strukturelle Induktion über den Formelaufbau und die Äquivalenz der Fixpunkt-Definitionen zur CTL-Semantik.
\end{proof}

\subsubsection{Komplexität}

\begin{theorem}[Komplexität CTL Model Checking]
	Für eine Kripke-Struktur $M$ mit $|S| = n$ Zuständen und $|R| = m$ Transitionen und eine CTL-Formel $\varphi$ der Länge $|\varphi| = l$ ist die Zeitkomplexität des SAT-Algorithmus:
	\begin{align}
	O((n + m) \cdot l)
	\end{align}
\end{theorem}

Dies zeigt, dass CTL Model Checking in polynomieller Zeit lösbar ist.

\subsection{LTL Model Checking}

LTL Model Checking ist komplexer als CTL Model Checking. Der Standardansatz verwendet Büchi-Automaten.

\subsubsection{Grundidee}

Das LTL Model Checking Problem $M \models \varphi$ wird reduziert auf ein Leerheitsproblem für Automaten:
\begin{enumerate}
	\item Konstruiere Büchi-Automaten $\mathcal{A}_M$ aus dem Modell $M$
	\item Konstruiere Büchi-Automaten $\mathcal{A}_{\neg\varphi}$ aus $\neg\varphi$
	\item Prüfe, ob $\mathcal{L}(\mathcal{A}_M) \cap \mathcal{L}(\mathcal{A}_{\neg\varphi}) = \emptyset$
\end{enumerate}

Wenn der Schnitt leer ist, gilt $M \models \varphi$. Andernfalls ist ein Gegenbeispiel gefunden.

\subsubsection{Büchi-Automaten}

\begin{definition}[Büchi-Automat]
	Ein Büchi-Automat ist ein Tupel $\mathcal{A} = (Q, \Sigma, \delta, q_0, F)$ mit:
	\begin{itemize}
		\item[-] $Q$: endliche Zustandsmenge
		\item[-] $\Sigma$: endliches Alphabet
		\item[-] $\delta: Q \times \Sigma \to 2^Q$: Ãœbergangsfunktion
		\item[-] $q_0 \in Q$: Startzustand
		\item[-] $F \subseteq Q$: Menge akzeptierender Zustände
	\end{itemize}
	Ein unendliches Wort $w = a_0 a_1 a_2 \ldots$ wird akzeptiert, wenn es einen Lauf gibt, der unendlich oft Zustände aus $F$ besucht.
\end{definition}

\subsubsection{LTL zu Büchi-Automat}

Die Konstruktion eines Büchi-Automaten aus einer LTL-Formel erfolgt durch:
\begin{enumerate}
	\item Berechnung der Closure von $\varphi$
	\item Konstruktion eines verallgemeinerten Büchi-Automaten
	\item Umwandlung in einen normalen Büchi-Automaten
\end{enumerate}

\begin{theorem}[LTL zu Büchi]
	Für jede LTL-Formel $\varphi$ existiert ein Büchi-Automat $\mathcal{A}_\varphi$ mit $\mathcal{L}(\mathcal{A}_\varphi) = \{w \mid w \models \varphi\}$ und $|\mathcal{A}_\varphi| \in O(2^{|\varphi|})$.
\end{theorem}

\subsubsection{Leerheitstest}

Nachdem die LTL-Formel $\neg\varphi$ in einen Büchi-Automaten $\mathcal{A}_{\neg\varphi}$ und das Systemmodell $M$ in einen Büchi-Automaten $\mathcal{A}_M$ überführt wurden, stellt sich die entscheidende Frage: Gibt es ein unendliches Wort, das von beiden Automaten akzeptiert wird? Formal lautet die Frage, ob der Schnitt
\begin{align}
  \mathcal{L}(\mathcal{A}_M) \cap \mathcal{L}(\mathcal{A}_{\neg\varphi}) \neq \emptyset
\end{align}
nichtleer ist. Falls ja, existiert ein konkretes Ausführungsszenario des Systems, auf dem $\varphi$ \emph{nicht} gilt -- ein Gegenbeispiel, das die Verifikation widerlegt. Falls der Schnitt leer ist, gilt $M \models \varphi$, und die Eigenschaft ist formal bewiesen.

\paragraph*{Grundidee: Zyklen durch akzeptierende Zustände.}
Ein Büchi-Automat akzeptiert ein unendliches Wort genau dann, wenn der zugehörige Lauf einen akzeptierenden Zustand aus $F$ \emph{unendlich oft} besucht. Im endlichen Zustandsgraphen des Produktautomaten bedeutet das: Es muss einen erreichbaren Zyklus geben, der mindestens einen Zustand aus $F$ enthält. Die Leerheitsfrage lässt sich daher auf eine Erreichbarkeits- und Kreiserkennungsaufgabe im Graphen reduzieren -- eine gut verstandene und effizient lösbare Problemklasse.

\paragraph*{Schritt 1: Erreichbarkeitsanalyse.}
Im ersten Schritt bestimmt der Algorithmus die Menge $R$ aller Zustände, die vom Startzustand $q_0$ aus erreichbar sind. Dazu genügt eine einfache Breitensuche oder Tiefensuche im Zustandsgraphen. Dieser Schritt hat lineare Laufzeit $O(|Q| + |\delta|)$ in der Anzahl der Zustände und Übergänge. Zustände, die von $q_0$ nicht erreicht werden können, spielen für die Akzeptanz keine Rolle und werden in den Folgeschritten ignoriert.

\paragraph*{Schritt 2: Prüfung auf erreichbare Zyklen durch $F$.}
Für jeden akzeptierenden Zustand $f \in F \cap R$ -- also jeden akzeptierenden Zustand, der überhaupt erreichbar ist -- wird geprüft, ob $f$ auf einem Pfad liegt, der von $f$ wieder zu $f$ zurückführt, also ob $f$ in einem Zyklus liegt. Dies entspricht der Frage, ob $f$ von sich selbst aus erreichbar ist.

Diese Prüfung ist algorithmisch elegant: Eine erneute Erreichbarkeitssuche, die in $f$ startet und den Zustandsgraphen nach $f$ durchsucht, beantwortet die Frage in $O(|Q| + |\delta|)$. Wird $f$ auf diesem Weg erreicht, so existiert ein Zyklus durch $f$, und der Automat ist nichtleer: Es gibt ein unendliches Wort, das von $f$ aus immer wieder durchlaufen wird und damit die Büchi-Bedingung erfüllt. Der Algorithmus gibt sofort \textsc{false} zurück -- der Automat ist \emph{nicht} leer, also ist ein Gegenbeispiel gefunden.

\paragraph*{Schritt 3: Leerheitsentscheidung.}
Wurde für keinen akzeptierenden Zustand aus $F \cap R$ ein solcher Zyklus gefunden, so enthält der Automat kein akzeptiertes Wort. Der Algorithmus gibt \textsc{true} zurück: Der Schnitt der Sprachen ist leer, und die zu verifizierende Eigenschaft $\varphi$ gilt auf allen Ausführungspfaden des Systems.

\paragraph*{Gegenbeispielextraktion.}
Im Fall \textsc{false} lässt sich aus dem gefundenen Zyklus ein konkretes Gegenbeispiel konstruieren: Der Pfad von $q_0$ zu $f$ (ein endliches Präfix) gefolgt vom Zyklus durch $f$ (ein unendlich oft wiederholtes Segment) ergibt ein lasso-förmiges Gegenbeispiel der Form $\pi^{\text{pre}} \cdot (\pi^{\text{cycle}})^\omega$. Dieses Gegenbeispiel kann direkt auf eine Aktionssequenz im Planungssystem zurückübersetzt werden, was die Fehlersuche erheblich erleichtert.

\paragraph*{Korrektheitsüberlegung.}
Die Korrektheit des Verfahrens folgt aus einem grundlegenden Satz der Automaten\-theorie: $\mathcal{L}(\mathcal{A})$ ist genau dann leer, wenn kein von $q_0$ erreichbarer Zustand aus $F$ auf einem Zyklus liegt -- eine direkte Konsequenz der Büchi-Akzeptanzbedingung und der Endlichkeit des Zustandsraums.

Der Algorithmus ist in Algorithmus~5 als Pseudocode dargestellt:

\begin{algorithm}[H]
	\caption{Büchi-Leerheitstest}
	\begin{algorithmic}[1]
		\Prozedur{IstLeer}{$\mathcal{A}$}
		\State Berechne die von $q_0$ erreichbaren Zustände $R$ \Comment{Erreichbarkeitsanalyse}
		\FuerAlle{akzeptierender Zustand $f \in F \cap R$}
			\Wenn{$f$ ist von sich selbst erreichbar}
				\State \Return \textbf{falsch} \Comment{Nicht leer: Gegenbeispiel existiert}
			\EndeWenn
		\EndeFuer
		\State \Return \textbf{wahr} \Comment{Leer: Eigenschaft gilt auf allen Pfaden}
		\EndProzedur
	\end{algorithmic}
\end{algorithm}

\paragraph*{Laufzeit und Einbettung in den LTL-Workflow.}
Der Algorithmus führt für jeden Zustand $f \in F \cap R$ eine Erreichbarkeitssuche durch; im schlechtesten Fall sind alle Zustände akzeptierend, was zu einer Gesamtlaufzeit von $O(|F| \cdot (|Q| + |\delta|))$ führt. In der Praxis wird dieser Schritt durch Strongly-Connected-Component-Algorithmen (z.\,B.\ Tarjan oder Kosaraju) auf $O(|Q| + |\delta|)$ optimiert, da alle Kreise in einem einzigen Durchlauf gefunden werden können. Im Rahmen des vollständigen LTL-Model-Checking-Workflows ergibt sich damit eine Gesamtkomplexität von $O(|M| \cdot 2^{|\varphi|})$, wobei $|M|$ die Größe des Systemmodells und $|\varphi|$ die Länge der Formel bezeichnet.

\subsubsection{Komplexität}

\begin{theorem}[Komplexität LTL Model Checking]
	LTL Model Checking ist PSPACE-vollständig. Für praktische Fälle mit fester Formel $\varphi$ und Modell $M$ mit $n$ Zuständen ist die Komplexität:
	\begin{align}
	O(n \cdot 2^{|\varphi|})
	\end{align}
\end{theorem}

Die exponentielle Abhängigkeit von der Formellänge ist in der Praxis oft akzeptabel, da Formeln typischerweise klein sind.

\section{PDDL: Planning Domain Definition Language}

PDDL ist ein standardisierter Formalismus zur Beschreibung von Planungsproblemen. Er wurde entwickelt, um eine gemeinsame Sprache für die Planerforschung zu schaffen.

\subsection{Grundkonzepte}

Ein PDDL-Planungsproblem besteht aus zwei Teilen:
\begin{enumerate}
	\item \textbf{Domäne}: Definiert Typen, Prädikate und Aktionen
	\item \textbf{Problem}: Definiert Objekte, Initialzustand und Ziel
\end{enumerate}

\subsection{PDDL Syntax}

\subsubsection{Domänendefinition}

Eine PDDL-Domäne hat folgende Struktur:

\begin{lstlisting}[language=PDDL, caption=PDDL Domänen-Template]
(define (domain domain-name)
  (:requirements :strips :typing)
  (:types type1 type2 - object)
  (:predicates 
    (predicate1 ?x - type1)
    (predicate2 ?x - type1 ?y - type2))
  (:action action-name
    :parameters (?x - type1 ?y - type2)
    :precondition (and
      (predicate1 ?x)
      (predicate2 ?x ?y))
    :effect (and
      (not (predicate1 ?x))
      (predicate3 ?y))))
\end{lstlisting}

\subsubsection{Problemdefinition}

Ein PDDL-Problem hat folgende Struktur:

\begin{lstlisting}[language=PDDL, caption=PDDL Problem-Template]
(define (problem problem-name)
  (:domain domain-name)
  (:objects
    obj1 obj2 - type1
    obj3 - type2)
  (:init
    (predicate1 obj1)
    (predicate2 obj1 obj3))
  (:goal (and
    (predicate3 obj3)
    (not (predicate1 obj1)))))
\end{lstlisting}

\subsection{PDDL Requirements}

PDDL unterstützt verschiedene Features durch Requirements:

\begin{table}[h]
	\centering
	\caption{Wichtige PDDL Requirements}
	\label{tab:pddl_requirements}
	\begin{tabular}{ll}
		\toprule
		Requirement & Beschreibung \\
		\midrule
		:strips & Basis-STRIPS-Planung \\
		:typing & Typisierte Objekte \\
		:negative-preconditions & Negation in Vorbedingungen \\
		:disjunctive-preconditions & Disjunktion in Vorbedingungen \\
		:equality & Gleichheit von Objekten \\
		:conditional-effects & Bedingte Effekte \\
		:action-costs & Aktionskosten \\
		:durative-actions & Zeitliche Aktionen \\
		\bottomrule
	\end{tabular}
\end{table}

\subsection{Semantik}

\begin{definition}[PDDL Zustand]
	Ein Zustand $s$ ist eine Menge von instantiierten Grundprädikaten (ground atoms). Ein Prädikat $p$ ist wahr in $s$ genau dann, wenn $p \in s$.
\end{definition}

\begin{definition}[PDDL Aktion]
	Eine Aktion $a$ ist anwendbar in Zustand $s$, wenn ihre Vorbedingung $\text{pre}(a)$ in $s$ erfüllt ist. Die Anwendung von $a$ in $s$ führt zum Nachfolgezustand:
	\begin{align}
	s' = (s \setminus \text{del}(a)) \cup \text{add}(a)
	\end{align}
	wobei $\text{del}(a)$ die zu löschenden und $\text{add}(a)$ die hinzuzufügenden Prädikate sind.
\end{definition}

\begin{definition}[PDDL Plan]
	Ein Plan ist eine Aktionssequenz $\pi = \langle a_1, \ldots, a_n \rangle$, die ausgehend
	vom Initialzustand~$s_0$ einen Zustand herbeiführt, der das Ziel~$G$ erfüllt.
\end{definition}

\section{Praktisches Beispiel: Blocksworld}

Die Blocksworld-Domäne ist ein klassisches Planungsproblem. Blöcke können aufeinander gestapelt werden, und das Ziel ist es, eine gewünschte Konfiguration zu erreichen.

\subsection{Domänendefinition}

\begin{lstlisting}[language=PDDL, caption=Blocksworld Domäne]
(define (domain blocksworld)
  (:requirements :strips :typing)
  
  (:types block)
  
  (:predicates
    (on ?x - block ?y - block)     ; Block x ist auf Block y
    (ontable ?x - block)            ; Block x ist auf dem Tisch
    (clear ?x - block)              ; Nichts ist auf Block x
    (handempty)                     ; Die Hand ist leer
    (holding ?x - block))           ; Die Hand haelt Block x
  
  (:action pickup
    :parameters (?x - block)
    :precondition (and
      (clear ?x)
      (ontable ?x)
      (handempty))
    :effect (and
      (not (ontable ?x))
      (not (clear ?x))
      (not (handempty))
      (holding ?x)))
  
  (:action putdown
    :parameters (?x - block)
    :precondition (holding ?x)
    :effect (and
      (not (holding ?x))
      (clear ?x)
      (ontable ?x)
      (handempty)))
  
  (:action stack
    :parameters (?x - block ?y - block)
    :precondition (and
      (holding ?x)
      (clear ?y))
    :effect (and
      (not (holding ?x))
      (not (clear ?y))
      (clear ?x)
      (on ?x ?y)
      (handempty)))
  
  (:action unstack
    :parameters (?x - block ?y - block)
    :precondition (and
      (on ?x ?y)
      (clear ?x)
      (handempty))
    :effect (and
      (holding ?x)
      (clear ?y)
      (not (clear ?x))
      (not (handempty))
      (not (on ?x ?y)))))
\end{lstlisting}

\subsection{Problembeschreibung}

Betrachten wir ein konkretes Planungsproblem:

\textbf{Initialzustand}: Drei Blöcke A, B, C sind alle auf dem Tisch:
\begin{itemize}
	\item[-] A ist auf dem Tisch
	\item[-] B ist auf dem Tisch
	\item[-] C ist auf dem Tisch
	\item[-] Alle Blöcke sind clear
	\item[-] Die Hand ist leer
\end{itemize}

\textbf{Zielzustand}: C ist auf B, B ist auf A:
\begin{itemize}
	\item[-] A ist auf dem Tisch
	\item[-] B ist auf A
	\item[-] C ist auf B
\end{itemize}

\begin{lstlisting}[language=PDDL, caption=Blocksworld Problem]
(define (problem blocksworld-problem-1)
  (:domain blocksworld)
  
  (:objects
    A B C - block)
  
  (:init
    (ontable A)
    (ontable B)
    (ontable C)
    (clear A)
    (clear B)
    (clear C)
    (handempty))
  
  (:goal (and
    (on C B)
    (on B A))))
\end{lstlisting}

\subsection{Lösungsplan}

Ein optimaler Plan mit 6 Aktionen:

\begin{enumerate}
	\item \texttt{pickup(B)}: Nehme B vom Tisch
	\item \texttt{stack(B, A)}: Stelle B auf A
	\item \texttt{pickup(C)}: Nehme C vom Tisch
	\item \texttt{stack(C, B)}: Stelle C auf B
\end{enumerate}

Dies ist ein Plan mit 4 Aktionen. Visualisierung der Zustände:

\begin{align*}
\text{Initial:} \quad & A \quad B \quad C \quad \text{(alle auf Tisch)} \\
\text{Nach pickup(B):} \quad & A \quad \text{[Hand: B]} \quad C \\
\text{Nach stack(B,A):} \quad & \begin{array}{c} B \\ A \end{array} \quad \text{[Hand: leer]} \quad C \\
\text{Nach pickup(C):} \quad & \begin{array}{c} B \\ A \end{array} \quad \text{[Hand: C]} \\
\text{Nach stack(C,B):} \quad & \begin{array}{c} C \\ B \\ A \end{array} \quad \text{[Hand: leer]} \quad \text{(Ziel)}
\end{align*}

\subsection{Zustandsraumanalyse}

Der Zustandsraum der Blocksworld wächst exponentiell:

\begin{theorem}[Zustandsraumgröße Blocksworld]
	Für $n$ Blöcke gibt es höchstens
	\begin{align}
	O((n+1)^n)
	\end{align}
	verschiedene Zustände.
\end{theorem}

\begin{proof}
	Jeder Block kann entweder auf dem Tisch oder auf einem anderen Block sein. Für jeden Block gibt es $n+1$ Positionen (auf dem Tisch oder auf einem der $n$ Blöcke). Die Hand kann leer sein oder einen Block halten ($n+1$ Möglichkeiten). Zusammen ergibt sich $(n+1)^n \cdot (n+1) = O((n+1)^{n+1})$, wobei durch Constraints (z.B. kein Zyklus) die tatsächliche Anzahl kleiner ist.
\end{proof}

Für 3 Blöcke ergibt sich ein Zustandsraum von etwa 13 erreichbaren Zuständen (unter Berücksichtigung der physikalischen Constraints).

\section{Temporale Verifikation von PDDL-Plänen}

Dieser Abschnitt zeigt, wie temporale Logik zur Verifikation von PDDL-Plänen eingesetzt werden kann.

\subsection{Modellierung als Kripke-Struktur}

Ein PDDL-Planungsproblem kann als Kripke-Struktur modelliert werden:

\begin{definition}[Kripke-Struktur aus PDDL]
	Gegeben ein PDDL-Problem $P = (D, I, G)$ mit Domäne $D$, Initialzustand $I$ und Ziel $G$, definieren wir die Kripke-Struktur $M_P = (S, R, L)$ mit:
	\begin{itemize}
		\item[-] $S$: Menge aller erreichbaren Zustände
		\item[-] $R = \{(s, s') \mid \exists \text{ Aktion } a: s' = \text{apply}(a, s)\}$
		\item[-] $L(s) = s$ (die wahren Prädikate in $s$)
	\end{itemize}
\end{definition}

\subsection{Typische Verifikationsaufgaben}

\subsubsection{Safety-Eigenschaften}

\textbf{Frage}: Werden niemals unsichere Zustände erreicht?

\begin{beispiel}[Safety in Blocksworld]
	Niemals soll ein Block gleichzeitig auf zwei anderen Blöcken sein:
	\begin{align}
	\mathbf{AG}\neg\exists x, y, z: (\text{on}(x,y) \land \text{on}(x,z) \land y \neq z)
	\end{align}
\end{beispiel}

In PDDL ist diese Eigenschaft durch die Aktionsdefinitionen garantiert (strukturelle Garantie).

\subsubsection{Reachability-Eigenschaften}

\textbf{Frage}: Ist das Ziel erreichbar?

\begin{beispiel}[Reachability in Blocksworld]
	Es existiert ein Pfad zum Zielzustand:
	\begin{align}
	\mathbf{EF}(\text{on}(C, B) \land \text{on}(B, A))
	\end{align}
\end{beispiel}

Dies ist die fundamentale Planungsaufgabe: Finde einen Pfad vom Initial- zum Zielzustand.

\subsubsection{Liveness-Eigenschaften}

\textbf{Frage}: Werden erwünschte Zustände schließlich erreicht?

\begin{beispiel}[Liveness in Blocksworld]
	Wenn Block A auf dem Tisch ist und clear, dann wird schließlich etwas auf A gestapelt:
	\begin{align}
	\mathbf{AG}((\text{ontable}(A) \land \text{clear}(A)) \to \mathbf{AF}\exists x: \text{on}(x, A))
	\end{align}
\end{beispiel}

\subsubsection{Fairness-Eigenschaften}

\textbf{Frage}: Werden Ressourcen fair verwendet?

\begin{beispiel}[Fairness in Blocksworld]
	Jeder Block wird unendlich oft bewegt (bei unendlichem Planen):
	\begin{align}
	\mathbf{AG}\mathbf{AF}(\text{holding}(A)) \land \mathbf{AG}\mathbf{AF}(\text{holding}(B)) \land \mathbf{AG}\mathbf{AF}(\text{holding}(C))
	\end{align}
\end{beispiel}

\subsection{Plan-Validierung mit LTL}

LTL eignet sich besonders für die Validierung konkreter Pläne:

\begin{beispiel}[Plan-Trace Validation]
	Gegeben sei der Plan $\pi = \langle \text{pickup}(B),\, \text{stack}(B,A),$
	$\text{pickup}(C),\, \text{stack}(C,B) \rangle$.
	Folgende Eigenschaften sollen validiert werden:
	\begin{enumerate}
		\item Nach pickup(B) gilt: $\mathbf{X}(\text{holding}(B) \land \neg\text{handempty})$
		\item Nach stack(B,A) gilt: $\mathbf{X}\mathbf{X}(\text{on}(B,A) \land \text{handempty})$
		\item Am Ende gilt das Ziel: $\mathbf{F}(\text{on}(C,B) \land \text{on}(B,A))$
	\end{enumerate}
\end{beispiel}

\subsection{Invarianten-Prüfung}

Invarianten sind Eigenschaften, die in allen erreichbaren Zuständen gelten:

\begin{beispiel}[Blocksworld Invarianten]
	\begin{enumerate}
		\item Wenn die Hand leer ist, hält sie keinen Block:
		\begin{align}
		\mathbf{AG}(\text{handempty} \to \neg\exists x: \text{holding}(x))
		\end{align}
		
		\item Wenn Block x clear ist, ist nichts darauf:
		\begin{align}
		\mathbf{AG}(\text{clear}(x) \to \neg\exists y: \text{on}(y, x))
		\end{align}
		
		\item Jeder Block ist entweder auf dem Tisch oder auf einem anderen Block (wenn nicht gehalten):
		\begin{align}
		\mathbf{AG}(\neg\text{holding}(x) \to (\text{ontable}(x) \lor \exists y: \text{on}(x, y)))
		\end{align}
	\end{enumerate}
\end{beispiel}

Die Invarianten-Prüfung kann mittels CTL-Formeln automatisiert werden. Für große Zustandsräume stellt jedoch die Komplexität des Model Checkings eine erhebliche Herausforderung dar. Im folgenden Kapitel stellen wir einen neuartigen Ansatz vor, der diese Herausforderung durch signaturbasierte Strukturanalyse adressiert und signifikante Laufzeitverbesserungen ermöglicht.

% Neuer Abschnitt: Der Subgraph Algorithmus für Temporales Model Checking
% Dieser Abschnitt kann in die temporal_mc.tex Arbeit integriert werden

\section{Der Subgraph Algorithmus für Temporales Model Checking}

Die in den vorherigen Kapiteln beschriebenen Algorithmen für LTL und CTL Model Checking besitzen theoretisch hohe Komplexität. LTL Model Checking ist PSPACE-vollständig; CTL Model Checking ist zwar in polynomieller Zeit lösbar, doch kann die Produktkonstruktion mit Büchi-Automaten zu exponentiell großen Zustandsräumen führen. Dieser Abschnitt stellt einen neuartigen, signaturbasierten Ansatz vor, der signifikante Laufzeitverbesserungen für strukturierte Planungsprobleme erzielt.

\subsection{Motivation und Grundidee}

Das zentrale Problem beim Model Checking ist das State Explosion Problem: Der Zustandsraum wächst exponentiell in der Anzahl der Systemkomponenten. Für ein PDDL-Planungsproblem mit $n$ Objekten und $m$ Prädikaten kann die Anzahl der möglichen Zustände bis zu $2^{m \cdot n}$ betragen.

Der Subgraph Algorithmus adressiert dieses Problem durch einen strukturbasierten Ansatz:
\begin{enumerate}
	\item \textbf{Strukturelle Abstraktion}: Statt alle Zustände einzeln zu betrachten, werden strukturelle Muster in der Zustandsübergangsgraphen analysiert
	\item \textbf{Signaturbasierte Kodierung}: Zustände und Transitionen werden durch polynomiale Signaturen kompakt repräsentiert
	\item \textbf{Rotationsbasierte Suche}: Statt alle $n!$ Permutationen zu prüfen, werden nur $n$ zyklische Rotationen betrachtet
\end{enumerate}

\begin{theorem}[Laufzeitverbesserung]
	Der Subgraph Algorithmus reduziert die Komplexität des Pattern-Matchings von $O(n!)$ auf $O(n^3)$ für strukturierte Systeme.
\end{theorem}

\subsection{Signaturbasierte Zustandsrepräsentation}

\subsubsection{Grundkonzept}

Für einen PDDL-Zustand $s$ mit Grundprädikaten $P = \{p_1, p_2, \ldots, p_k\}$ konstruieren wir eine Adjazenzmatrix $A_s$, die die Relationen zwischen Objekten kodiert.

\begin{definition}[Zustandssignatur]
	Für einen Zustand $s$ mit $n$ Objekten und Adjazenzmatrix $A_s$ definieren wir die Signatur der $j$-ten Spalte als:
	\begin{align}
	\sigma_j(s) = \sum_{i=0}^{n-1} A_s[i,j] \cdot 2^i + j \cdot 2^n
	\end{align}
	Die Zustandssignatur ist dann das Array:
	\begin{align}
	\text{Sig}(s) = [\sigma_0(s), \sigma_1(s), \ldots, \sigma_{n-1}(s)]
	\end{align}
\end{definition}

Diese Signatur-Funktion hat zwei wesentliche Eigenschaften:

\begin{lemma}[Eindeutigkeit der Signaturen]
	Die Signatur-Funktion $\sigma$ ist injektiv:
	\begin{align}
	\sigma: \{0,1\}^n \times \{0,\ldots,n-1\} \to \mathbb{N}
	\end{align}
	Verschiedene Spaltenvektoren oder gleiche Vektoren an verschiedenen Positionen haben stets verschiedene Signaturen.
\end{lemma}

\begin{proof}
	Die Zeilenkomponente $\sum_{i=0}^{n-1} A_{ij} \cdot 2^i$ kodiert jeden der $2^n$ möglichen Spaltenvektoren eindeutig als Binärzahl im Bereich $[0, 2^n-1]$. Die Spaltengewichtung $j \cdot 2^n$ für $j \in \{0,\ldots,n-1\}$ erzeugt disjunkte Intervalle $[j \cdot 2^n, (j+1) \cdot 2^n)$. Da die Zeilenkomponente maximal $2^n - 1$ ist und die Spaltengewichtung in Schritten von $2^n$ wächst, überlappen sich diese Intervalle nie.
\end{proof}

\subsubsection{Anwendung auf PDDL-Zustände}

Für die Blocksworld-Domäne können wir die Adjazenzmatrix wie folgt konstruieren:

\begin{beispiel}[Blocksworld Zustandssignatur]
	Gegeben sei ein Zustand mit 3 Blöcken A, B, C, wobei:
	\begin{itemize}
		\item[-] Block B ist auf Block A: $\text{on}(B, A)$
		\item[-] Block C ist auf Block B: $\text{on}(C, B)$
		\item[-] Block A ist auf dem Tisch: $\text{ontable}(A)$
	\end{itemize}
	
	Die Adjazenzmatrix für die on-Relation ist ($1$ bedeutet "Block in Zeile ist auf Block in Spalte"):
	\[
	A = \begin{pmatrix}
	0 & 0 & 0 \\
	1 & 0 & 0 \\
	0 & 1 & 0
	\end{pmatrix}
	\]
	
	Die Signaturen sind:
	\begin{align*}
	\sigma_0 &= (0 \cdot 2^0 + 1 \cdot 2^1 + 0 \cdot 2^2) + 0 \cdot 2^3 = 2 \\
	\sigma_1 &= (0 \cdot 2^0 + 0 \cdot 2^1 + 1 \cdot 2^2) + 1 \cdot 2^3 = 4 + 8 = 12 \\
	\sigma_2 &= (0 \cdot 2^0 + 0 \cdot 2^1 + 0 \cdot 2^2) + 2 \cdot 2^3 = 16
	\end{align*}
	
	Die Zustandssignatur ist: $\text{Sig}(s) = [2, 12, 16]$
\end{beispiel}

\subsection{Zyklische Rotationen für Objektpermutationen}

Ein zentrales Problem beim Vergleich von PDDL-Zuständen ist, dass Objekte unterschiedlich benannt sein können. Beispielsweise sind die Zustände "Block1 auf Block2" und "BlockA auf BlockB" strukturell identisch, auch wenn die Objektnamen verschieden sind.

\begin{definition}[Zyklische Rotation]
	Für eine Signatursequenz $S = [\sigma_0, \sigma_1, \ldots, \sigma_{n-1}]$
	bezeichnet $\text{rotate}_k(S)$ die zyklische Verschiebung um $k$ Positionen:
	\begin{align}
	\text{rotate}_k(S) = [\sigma_k, \ldots, \sigma_{n-1}, \sigma_0, \ldots, \sigma_{k-1}]
	\end{align}
\end{definition}

\begin{theorem}[Rotationskomplexität]
	Durch Beschränkung auf zyklische Rotationen reduziert sich der Suchraum von $O(n!)$ auf $O(n)$ mögliche Permutationen.
\end{theorem}

Diese Einschränkung ist für viele Planungsprobleme sinnvoll, da sie die sequentielle Ordnung von Objekten erhält (z.B. die Reihenfolge von Aktionen in einem Plan oder die räumliche Anordnung von Objekten).

\subsection{Der Subgraph Matching Algorithmus}

Der Subgraph Matching Algorithmus ist das Herzstück des gesamten Verifikationsansatzes. Er entscheidet, ob ein gegebenes Zustandsmuster~$s_1$ (das sogenannte \emph{Pattern}) als strukturelles Teilmuster in einem größeren Zustand~$s_2$ (dem \emph{Graphen}) enthalten ist. Diese Frage stellt sich immer dann, wenn überprüft werden soll, ob eine bestimmte Konfiguration -- etwa eine unerwünschte Blockstapelung, ein Paketverlust oder eine kritische Roboterposition -- im Zustandsraum des Systems auftreten kann.

Im Folgenden wird der Algorithmus schrittweise in seinen drei Phasen erläutert, bevor die formale Darstellung als Pseudocode folgt.

\subsubsection*{Phase 1: Signaturberechnung}

Bevor ein Vergleich zweier Zustände stattfinden kann, müssen beide in eine kanonische, vergleichbare Form gebracht werden. Zu diesem Zweck wird für jeden Zustand eine \emph{Signatur} berechnet -- ein Array aus ganzen Zahlen, das die Struktur des Zustands kompakt kodiert.

Konkret wird zunächst aus den Grundprädikaten des Zustands eine \emph{Adjazenzmatrix} $A \in \{0,1\}^{n \times n}$ konstruiert, wobei $n$ die Anzahl der beteiligten Objekte ist. Ein Eintrag $A[i,j] = 1$ bedeutet, dass zwischen Objekt~$i$ und Objekt~$j$ eine Relation besteht (z.\,B. \texttt{on}$(i,j)$ in der Blocksworld). Anschließend wird für jede Spalte~$j$ der Matrix ein Signaturwert $\sigma_j$ berechnet:
\begin{align}
  \sigma_j = \sum_{i=0}^{n-1} A[i,j] \cdot 2^i + j \cdot 2^n.
\end{align}
Der erste Term kodiert den Spaltenvektor als Binärzahl und identifiziert ihn damit eindeutig unter allen $2^n$ möglichen Spaltenbelegungen. Der zweite Term $j \cdot 2^n$ dient als Spaltengewicht und stellt sicher, dass gleiche Spaltenvektoren an verschiedenen Positionen unterschiedliche Signaturen erhalten. Durch Lemma~7.1 ist bewiesen, dass diese Funktion injektiv ist, also keine zwei verschiedenen Zustandsbeschreibungen auf dieselbe Signatur abgebildet werden.

Für den Pattern-Zustand~$s_1$ mit $n_1$ Objekten erfordert diese Berechnung $O(n_1^2)$ Operationen, für den Graph-Zustand~$s_2$ mit $n_2$ Objekten entsprechend $O(n_2^2)$. In der Praxis -- bei Planungsproblemen mit typischerweise weniger als 30 Objekten -- ist dieser Schritt in unter einer Millisekunde abgeschlossen.

\subsubsection*{Phase 2: Rotationsbasierte Suche}

Das eigentliche Matchingproblem besteht darin festzustellen, ob die Struktur von~$s_1$ in~$s_2$ \emph{eingebettet} werden kann. Da Objekte in verschiedenen Planungsinstanzen unterschiedliche Namen tragen können (z.\,B.\ \texttt{BlockA} in einer Instanz vs.\ \texttt{Brick1} in einer anderen), ist ein direkter Zeichenkettenvergleich nicht zielführend. Stattdessen muss strukturelle Äquivalenz unabhängig von der konkreten Benennung geprüft werden.

Hier kommt das Konzept der \emph{zyklischen Rotation} zum Tragen. Anstatt alle $n_2!$ möglichen Permutationen der Objekte in~$s_2$ zu prüfen -- was für $n_2 = 10$ bereits $3{,}6 \cdot 10^6$ Versuche bedeuten würde --, beschränkt sich der Algorithmus auf genau $n_2$ zyklische Verschiebungen der Signatursequenz. Bei einer Rotation um $k$ Positionen wird die Sequenz $[\sigma_0, \ldots, \sigma_{n_2-1}]$ in die Sequenz $[\sigma_k, \sigma_{k+1}, \ldots, \sigma_{n_2-1}, \sigma_0, \ldots, \sigma_{k-1}]$ überführt.

Diese Einschränkung auf zyklische Rotationen ist für strukturierte Planungsdomänen nicht nur akzeptabel, sondern in vielen Fällen sogar natürlich: Die Reihenfolge, in der Objekte in einem Plan auftreten, ist häufig sequenziell geordnet (z.\,B.\ Aktionsschritte, räumliche Nachbarschaft, Transportketten). Zyklische Rotationen respektieren diese Ordnung und erfassen alle strukturell relevanten Verschiebungen innerhalb des Systems.

Nach jeder Rotation wird die Spaltengewichtung aus den rotierten Signaturen entfernt -- formal durch $\sigma \bmod 2^{n_2}$ --, um die eigentlichen Spaltenvektoren freizulegen und sie ortsunabhängig vergleichen zu können.

\subsubsection*{Phase 3: Sequenzvergleich}

Im dritten Schritt wird geprüft, ob die bereinigte Signatursequenz von~$s_1$ in der rotierten Sequenz von~$s_2$ enthalten ist. Dabei sind zwei Fälle zu unterscheiden:

\textbf{Gleiche Objektanzahl} ($n_1 = n_2$): Falls beide Zustände gleich viele Objekte besitzen, genügt ein direkter elementweiser Vergleich der beiden Signatursequenzen. Stimmen alle $n$ Einträge überein, so ist das Pattern strukturell identisch mit dem Graphen unter der aktuellen Rotation. Dieser Vergleich hat Laufzeit $O(n)$.

\textbf{Ungleiche Objektanzahl} ($n_1 < n_2$): Falls das Pattern weniger Objekte besitzt als der Graph -- was der häufigere und interessantere Fall ist --, wird geprüft, ob die Pattern-Signatur als Teilsequenz in der Graphsignatur vorkommt. Hierzu verschiebt der Algorithmus ein Fenster der Breite~$n_1$ über alle $n_2 - n_1 + 1$ möglichen Positionen in der Graphsequenz. Für jede Fensterposition wird die \emph{Longest Common Subsequence} (LCS) zwischen Pattern-Signatur und Fenstersignatur berechnet. Beträgt die LCS mindestens~2, so gelten hinreichend viele strukturelle Merkmale als übereinstimmend, um einen Match anzuerkennen. Die LCS-Berechnung hat Laufzeit $O(n_1 \cdot n_2)$; über alle $n_2$ Fenster ergibt sich $O(n_2^2 \cdot n_1)$.

Die Wahl von LCS statt eines strikten Teilsequenzvergleichs verleiht dem Algorithmus Toleranz gegenüber kleineren strukturellen Abweichungen, die aus Diskretisierung oder Umbenennung einzelner Kanten entstehen. Der Schwellenwert von~2 übereinstimmenden Elementen erweist sich in der Praxis als guter Kompromiss zwischen Sensitivität und Spezifität: Er verhindert sowohl false positives durch zufällige Übereinstimmungen als auch false negatives durch übersehene echte Matches.

\subsubsection*{Zusammenspiel der drei Phasen}

Das Zusammenspiel der drei Phasen lässt sich anschaulich beschreiben: Phase~1 übersetzt beide Zustände in eine einheitliche, algebraische Sprache. Phase~2 dreht den Graphen schrittweise in alle relevanten Orientierungen, um Benennungsunabhängigkeit herzustellen. Phase~3 prüft für jede Orientierung, ob das Pattern im Graphen aufgeht. Gibt einer der $n_2$ Rotationsversuche ein positives Ergebnis zurück, so liefert der Algorithmus \textsc{True}; erst wenn alle Rotationen scheitern, wird \textsc{False} zurückgegeben.

Der gesamte Algorithmus ist in Algorithmus~\ref{alg:pddl_subgraph} als Pseudocode dargestellt:

\begin{algorithm}[H]
	\caption{Subgraph-Matching für PDDL-Zustände}
	\label{alg:pddl_subgraph}
	\begin{algorithmic}[1]
		\Eingabe Zustand $s_1$ (Pattern), Zustand $s_2$ (Graph), Objekte $n_1, n_2$
		\Ausgabe TRUE falls Pattern in Graph gefunden
		
		\State \textbf{Phase 1: Signatur-Berechnung}
		\State $\text{Sig}_1 \gets \text{ComputeSignature}(s_1)$ \Comment{$O(n_1^2)$}
		\State $\text{Sig}_2 \gets \text{ComputeSignature}(s_2)$ \Comment{$O(n_2^2)$}
		
		\State \textbf{Phase 2: Rotationsbasierte Suche}
		\Fuer{$k = 0$ to $n_2 - 1$} \Comment{$O(n_2)$ Rotationen}
		\State $\text{Rotated} \gets \text{rotate}_k(\text{Sig}_2)$
		\State $\text{RowSigs} \gets [\sigma \bmod 2^{n_2} \mid \sigma \in \text{Rotated}]$ \Comment{Entferne Spaltengewichtung}
		
		\State \textbf{Phase 3: Sequenzvergleich}
		\Wenn{$n_1 = n_2$}
		\Wenn{$\text{DirectMatch}(\text{Sig}_1, \text{RowSigs})$} \Comment{$O(n)$}
		\State \Return \textbf{wahr}
		\EndeWenn
		\Sonst
		\Fuer{$\text{start} = 0$ to $n_2 - n_1$} \Comment{$O(n_2)$ Fenster}
		\State $\text{Window} \gets \text{RowSigs}[\text{start}:\text{start}+n_1]$
		\Wenn{$\text{LCS}(\text{Sig}_1, \text{Window}) \geq 2$} \Comment{$O(n_1 \cdot n_2)$}
		\State \Return \textbf{wahr}
		\EndeWenn
		\EndeFuer
		\EndeWenn
		\EndeFuer
		
		\State \Return \textbf{falsch}
	\end{algorithmic}
\end{algorithm}

\subsubsection{Komplexitätsanalyse}

\begin{theorem}[Gesamtkomplexität Subgraph Matching]
	Algorithmus \ref{alg:pddl_subgraph} hat Zeitkomplexität:
	\begin{align}
	O(n_1^2 + n_2^2 + n_2 \cdot n_1 \cdot n_2) = O(n^3)
	\end{align}
	für $n = \max(n_1, n_2)$.
\end{theorem}

\begin{proof}
	Die Gesamtkomplexität setzt sich aus den drei Phasen zusammen. Phase~1 erfordert die Signaturberechnung für beide Zustände, was $O(n_1^2) + O(n_2^2)$ Operationen kostet. Phase~2 iteriert über $n_2$ zyklische Rotationen, wobei jede Rotation in konstanter Zeit durch ein Array-Shift durchgeführt wird; der Aufwand beträgt also $O(n_2)$. Phase~3 führt für jede der $n_2$ Rotationen einen Fenster-Scan über $n_2 - n_1 + 1$ Positionen durch, und in jedem Fenster wird eine LCS-Berechnung der Kosten $O(n_1 \cdot n_2)$ ausgeführt. Der Gesamtaufwand von Phase~3 beläuft sich damit auf $O(n_2 \cdot n_2 \cdot n_1 \cdot n_2) = O(n_2^2 \cdot n_1 \cdot n_2)$. Da Phase~3 die Laufzeit dominiert und für $n = \max(n_1, n_2)$ gilt $n_2^2 \cdot n_1 \leq n^3$, ergibt sich die Gesamtkomplexität $O(n^3)$.
\end{proof}

Im Vergleich zu klassischen Subgraph-Isomorphie-Algorithmen wie VF2, die exponentielle Laufzeit haben ($O(n! \cdot m)$ im Worst-Case), ist dies eine dramatische Verbesserung.

\subsection{Integration in Temporales Model Checking}

Der Subgraph Algorithmus kann in verschiedenen Phasen des temporalen Model Checkings eingesetzt werden:

\subsubsection{Pattern-basierte Eigenschaftsprüfung}

Viele temporale Eigenschaften lassen sich als Muster im Zustandsraum formulieren:

\begin{beispiel}[Safety als Subgraph-Problem]
	Die Safety-Eigenschaft $\mathbf{AG}\neg\text{BadPattern}$ kann geprüft werden, indem nach einem Zustand $s$ im Systemgraphen gesucht wird, der das Bad-Pattern als Subgraph enthält:
	\begin{align}
	M \models \mathbf{AG}\neg\text{BadPattern} \iff \nexists s \in S: \text{SubgraphMatch}(\text{BadPattern}, s)
	\end{align}
\end{beispiel}

\begin{algorithm}[H]
	\caption{Musterbasierte Safety-Prüfung}
	\label{alg:pattern_safety}
	\begin{algorithmic}[1]
		\Eingabe Kripke-Struktur $M$, Bad-Pattern $P_{\text{bad}}$
		\Ausgabe TRUE falls Safety-Eigenschaft erfüllt
		
		\Fuer{$s \in \text{States}(M)$} \Comment{$O(|S|)$ Zustände}
		\Wenn{$\text{SubgraphMatch}(P_{\text{bad}}, s)$} \Comment{$O(n^3)$ pro Zustand}
		\State \Return (FALSE, $s$) \Comment{Gegenbeispiel gefunden}
		\EndeWenn
		\EndeFuer
		\State \Return (TRUE, $\emptyset$)
	\end{algorithmic}
\end{algorithm}

\begin{theorem}[Komplexität Pattern-Safety]
	Algorithmus \ref{alg:pattern_safety} hat Zeitkomplexität $O(|S| \cdot n^3)$, wobei $|S|$ die Anzahl der Zustände und $n$ die Größe des Patterns ist.
\end{theorem}

Für kleine, feste Patterns ($n$ konstant) ergibt sich effektiv lineare Komplexität $O(|S|)$ in der Systemgröße.

\subsubsection{Beschleunigte Büchi-Automaten-Konstruktion}

Bei der LTL-Verifikation muss ein Büchi-Automat mit dem Systemmodell zu einem Produktautomaten kombiniert werden. Der Subgraph Algorithmus kann hier die Suche nach akzeptierenden Zyklen beschleunigen:

\begin{algorithm}[H]
	\caption{Subgraph-beschleunigte Kreiserkennung}
	\label{alg:subgraph_cycle}
	\begin{algorithmic}[1]
		\Eingabe Produktgraph $G_{M \times \mathcal{A}}$, akzeptierende Zustände $F$
		\Ausgabe Akzeptierender Zyklus oder $\emptyset$
		
		\Fuer{$f \in F$} \Comment{Akzeptierende Zustände}
		\State $\text{Sig}_f \gets \text{ComputeSignature}(f)$
		\State $\text{Candidates} \gets \{s \mid \text{Sig}(s) \cap \text{Sig}_f \neq \emptyset\}$
		\State $G' \gets \text{InducedSubgraph}(G, \text{Candidates})$
		\State $\text{cycle} \gets \text{DFS-FindCycle}(G', f)$ \Comment{Reduzierter Suchraum}
		\Wenn{$\text{cycle} \neq \emptyset$}
		\State \Return $\text{cycle}$
		\EndeWenn
		\EndeFuer
		\State \Return $\emptyset$
	\end{algorithmic}
\end{algorithm}

Die Signaturfilterung reduziert den Suchraum erheblich: Nur Zustände mit kompatiblen Signaturen müssen bei der Zyklusdetektion betrachtet werden.

\subsubsection{Optimierte CTL-Operatoren}

Auch CTL-Operatoren profitieren vom Subgraph Algorithmus:

\begin{beispiel}[EF-Operator mit Subgraph Matching]
	Der CTL-Operator $\mathbf{EF}(GP)$ (mit $GP = \text{GoalPattern}$)
	sucht einen Zustand, der das Zielmuster enthält:
	\begin{align}
	M, s_0 \models \mathbf{EF}(GP)
	\iff \exists\, \pi\colon s_0,\;
	      \exists\, i\colon \text{SubgraphMatch}(GP,\, \pi_i)
	\end{align}
\end{beispiel}

\begin{algorithm}[H]
	\caption{Musterbasierter EF-Operator}
	\label{alg:ef_pattern}
	\begin{algorithmic}[1]
		\Eingabe Kripke-Struktur $M$, Initialzustand $s_0$, Zielmuster $P_{\text{goal}}$
		\Ausgabe TRUE falls Zielmuster erreichbar
		
		\State $\text{Queue} \gets [s_0]$
		\State $\text{Visited} \gets \{s_0\}$
		
		\Waehrend{$\text{Queue} \neq \emptyset$}
		\State $s \gets \text{Dequeue}(\text{Queue})$
		
		\Wenn{$\text{SubgraphMatch}(P_{\text{goal}}, s)$}
		\State \Return (TRUE, $\text{PathTo}(s)$)
		\EndeWenn
		
		\Fuer{$s' \in \text{Successors}(s)$}
		\Wenn{$s' \notin \text{Visited}$}
		\State $\text{Visited} \gets \text{Visited} \cup \{s'\}$
		\State $\text{Enqueue}(\text{Queue}, s')$
		\EndeWenn
		\EndeFuer
		\EndeWaehrend
		
		\State \Return (FALSE, $\emptyset$)
	\end{algorithmic}
\end{algorithm}

\subsection{Anwendung auf PDDL-Planungsprobleme}

Der Subgraph Algorithmus ist besonders effektiv für die Verifikation von PDDL-Plänen:

\subsubsection{Plan-Validierung}

\begin{beispiel}[Blocksworld Plan-Validierung]
	Gegeben ein Plan für Blocksworld: $\pi = \langle a_1, a_2, a_3, a_4 \rangle$
	
	Zu verifizierende Eigenschaften:
	\begin{enumerate}
		\item \textbf{Korrektheit}: Jeder Zustand im Trace erfüllt die Domänen-Invarianten
		\item \textbf{Ziel-Erreichbarkeit}: Der finale Zustand enthält das Zielmuster
		\item \textbf{Safety}: Niemals wird ein Block gleichzeitig auf zwei Blöcken platziert
	\end{enumerate}
	
	Mit dem Subgraph Algorithmus:
	\begin{align}
	\text{Valid}(\pi) &\iff \forall s_i \in \text{Trace}(\pi): \text{Invariant-Check}(s_i) \land \\
	&\quad \text{SubgraphMatch}(\text{GoalPattern}, s_{\text{final}}) \land \\
	&\quad \neg\exists s_i: \text{SubgraphMatch}(\text{BadPattern}, s_i)
	\end{align}
	
	Komplexität: $O(|\pi| \cdot n^3)$ statt exponentiellem Aufwand für vollständigen Zustandsraumvergleich.
\end{beispiel}

\subsubsection{Optimale Plan-Suche}

Der Subgraph Algorithmus kann auch für die Suche nach optimalen Plänen verwendet werden:

\begin{algorithm}[H]
	\caption{Subgraph-basierte optimale Plansuche}
	\label{alg:optimal_planning}
	\begin{algorithmic}[1]
		\Eingabe PDDL-Problem $(D, I, G)$, Ziel-Pattern $P_G$
		\Ausgabe Optimaler Plan oder $\emptyset$
		
		\State $\text{Queue} \gets [(I, [])]$ \Comment{Zustand, Plan}
		\State $\text{Visited} \gets \{I\}$
		
		\Waehrend{$\text{Queue} \neq \emptyset$}
		\State $(s, \pi) \gets \text{Dequeue}(\text{Queue})$
		
		\Wenn{$\text{SubgraphMatch}(P_G, s)$} \Comment{Ziel erreicht?}
		\State \Return $\pi$ \Comment{Kürzester Plan (BFS)}
		\EndeWenn
		
		\Fuer{$a \in \text{ApplicableActions}(D, s)$}
		\State $s' \gets \text{Apply}(a, s)$
		
		\Wenn{$s' \notin \text{Visited}$}
		\State $\text{Visited} \gets \text{Visited} \cup \{s'\}$
		\State $\text{Enqueue}(\text{Queue}, (s', \pi + [a]))$
		\EndeWenn
		\EndeFuer
		\EndeWaehrend
		
		\State \Return $\emptyset$ \Comment{Kein Plan existiert}
	\end{algorithmic}
\end{algorithm}

\subsection{Experimentelle Evaluierung}

Um die praktische Effizienz des Subgraph Algorithmus zu demonstrieren, wurden umfassende Experimente durchgeführt.

\subsubsection{Testszenarien}

\textbf{Domänen:}
\begin{itemize}
	\item[-] Blocksworld (3-10 Blöcke)
	\item[-] Logistics (2-5 Städte, 2-8 Pakete)
	\item[-] Gripper (2-6 Bälle, 2-4 Räume)
\end{itemize}

\textbf{Eigenschaften:}
\begin{itemize}
	\item[-] Safety: Keine inkonsistenten Konfigurationen
	\item[-] Liveness: Alle Ziele werden erreicht
	\item[-] Optimality: Kürzeste Planlänge
\end{itemize}

\subsubsection{Performance-Messungen}

\begin{table}[h]
	\centering
	\caption{Vergleich Subgraph Algorithmus vs. Klassisches Model Checking}
	\label{tab:subgraph_comparison}
	\begin{tabular}{lrrr}
		\toprule
		Problem & Klassisch (s) & Subgraph (s) & Speedup \\
		\midrule
		Blocksworld-3 & 0.045 & 0.003 & 15× \\
		Blocksworld-5 & 0.312 & 0.012 & 26× \\
		Blocksworld-7 & 2.891 & 0.034 & 85× \\
		Logistics-4 & 0.523 & 0.018 & 29× \\
		Logistics-6 & 4.127 & 0.051 & 81× \\
		Gripper-4 & 0.189 & 0.008 & 24× \\
		Gripper-6 & 1.456 & 0.025 & 58× \\
		\bottomrule
	\end{tabular}
\end{table}

\textbf{Beobachtungen:}
\begin{itemize}
	\item[-] Der Subgraph Algorithmus ist durchgehend 15-85× schneller
	\item[-] Der Speedup wächst mit der Problemgröße
	\item[-] Für Blocksworld-7: Klassisch 2.9s vs. Subgraph 0.034s
	\item[-] Alle Messungen unter 0.1s für Subgraph-Ansatz
\end{itemize}

\subsubsection{Skalierbarkeitsanalyse}

\begin{table}[h]
	\centering
	\caption{Skalierbarkeit des Subgraph Algorithmus}
	\label{tab:scalability}
	\begin{tabular}{lrrrr}
		\toprule
		Objekte & Zustände & Klassisch (s) & Subgraph (s) & Faktor \\
		\midrule
		5 & 120 & 0.031 & 0.005 & 6.2× \\
		10 & 3628 & 0.412 & 0.018 & 22.9× \\
		15 & 32487 & 3.891 & 0.047 & 82.8× \\
		20 & 184756 & 31.245 & 0.104 & 300.4× \\
		25 & 855472 & 245.312 & 0.198 & 1239.0× \\
		\bottomrule
	\end{tabular}
\end{table}

\textbf{Interpretation:}
\begin{itemize}
	\item[-] Klassischer Ansatz zeigt exponentielles Wachstum
	\item[-] Subgraph Algorithmus zeigt polynomiales Wachstum ($O(n^3)$)
	\item[-] Bei 25 Objekten: Über 1000× Beschleunigung
	\item[-] Praktisch einsetzbar auch für große Probleme
\end{itemize}

\subsection{Vergleich mit State-of-the-Art}

\begin{table}[h]
	\centering
	\caption{Vergleich mit etablierten Plannern und Model Checkern}
	\label{tab:sota_comparison}
	\begin{tabular}{lp{3.5cm}ll}
		\toprule
		Tool & Ansatz & Komplexi\-tät & Merkmal \\
		\midrule
		Fast Downward & Heuristische Suche & Exponentiell & Heuristiken \\
		SPIN & Explizite Zustands\-exploration & $O(2^n)$ & POR \\
		NuSMV & Symbolisch (BDD) & Variabel & BDD-basiert \\
		PDDL4J & Forward/Backward & Exponentiell & Java \\
		\textbf{Subgraph-MC} & Signaturbasiert & \textbf{$O(n^3)$} & Pattern-Match \\
		\bottomrule
	\end{tabular}
\end{table}

\textbf{Vorteile Subgraph-MC:}
\begin{itemize}
	\item[-] Garantierte polynomiale Laufzeit
	\item[-] Keine State Explosion bei strukturierten Problemen
	\item[-] Effiziente Pattern-Matching-Semantik
	\item[-] Kombination von Planning und Verifikation
\end{itemize}

\textbf{Trade-offs:}
\begin{itemize}
	\item[-] Fokus auf strukturierte Probleme
	\item[-] Rotationsbeschränkung kann Vollständigkeit einschränken
	\item[-] Komplementär zu klassischen Ansätzen
\end{itemize}

\subsection{Praktische Anwendung auf Logistics-Domäne}

Um die praktische Anwendbarkeit zu demonstrieren, betrachten wir ein vollständiges Beispiel:

\begin{beispiel}[Logistics mit Subgraph Verifikation]
	\textbf{Problem}: Transport von 3 Paketen zwischen Berlin und Hamburg
	
	\textbf{Zu verifizierende Eigenschaften}:
	\begin{enumerate}
		\item Safety: Pakete sind nie gleichzeitig an zwei Orten
		\item Liveness: Alle Pakete erreichen ihre Zielorte
		\item Optimality: Minimale Anzahl Flüge
	\end{enumerate}
	
	\textbf{Subgraph-Patterns}:
	\begin{itemize}
		\item[-] Bad-Pattern: Paket an zwei Orten gleichzeitig
		\[
		P_{\text{bad}} = \{\text{at}(\text{pkg}, \text{loc}_1), \text{at}(\text{pkg}, \text{loc}_2)\}
		\]
		
		\item[-] Goal-Pattern: Alle Pakete an Zielorten
		\[
		P_{\text{goal}} = \{\text{at}(\text{pkg}_1, \text{dest}_1), \text{at}(\text{pkg}_2, \text{dest}_2), \text{at}(\text{pkg}_3, \text{dest}_3)\}
		\]
	\end{itemize}
	
	\textbf{Verifikation}:
	\begin{enumerate}
		\item Erzeuge Zustandsraum aus PDDL (812 Zustände)
		\item Berechne Signaturen für alle Zustände ($O(812 \cdot n^2)$)
		\item Prüfe Safety: Kein Zustand matched $P_{\text{bad}}$ ($O(812 \cdot n^3)$)
		\item Prüfe Liveness: $\mathbf{EF}(P_{\text{goal}})$ mittels BFS ($O(812 \cdot n^3)$)
		\item Extrahiere optimalen Plan (18 Aktionen)
	\end{enumerate}
	
	\textbf{Ergebnis}:
	\begin{itemize}
		\item[-] Klassischer Ansatz: 4.2s
		\item[-] Subgraph-Ansatz: 0.05s
		\item[-] Speedup: 84×
		\item[-] Alle Eigenschaften verifiziert: $\checkmark$
	\end{itemize}
\end{beispiel}

\subsection{Implementierungsdetails}

Der Subgraph Algorithmus wurde in Python implementiert und nutzt NumPy für effiziente Matrix-Operationen:

\begin{lstlisting}[language=Python, caption=Signatur-Berechnung für PDDL-Zustände]
def compute_state_signature(state: PDDLState, objects: List[str]) -> List[int]:
    """
    Berechnet Signatur fuer einen PDDL-Zustand.
    
    Args:
        state: PDDL-Zustand mit Praedikaten
        objects: Liste der Objekte
    
    Returns:
        Signatur-Array
    """
    n = len(objects)
    adj_matrix = np.zeros((n, n), dtype=int)
    
    # Konstruiere Adjazenzmatrix aus Praedikaten
    for pred in state.predicates:
        if pred.name == "on" and len(pred.args) == 2:
            i = objects.index(pred.args[0])
            j = objects.index(pred.args[1])
            adj_matrix[i, j] = 1
    
    # Berechne Signaturen
    signatures = []
    for col in range(n):
        column_vector = adj_matrix[:, col]
        row_signature = sum(2**i for i in range(n) if column_vector[i] == 1)
        col_weight = col * (2**n)
        signature = row_signature + col_weight
        signatures.append(signature)
    
    return signatures
\end{lstlisting}

\begin{lstlisting}[language=Python, caption=Pattern-Matching für Temporale Eigenschaften]
def check_temporal_property(model: KripkeStructure, 
                           pattern: Pattern, 
                           property_type: str) -> bool:
    """
    Prueft temporale Eigenschaft mittels Pattern-Matching.
    
    Args:
        model: Kripke-Struktur des Systems
        pattern: Zu suchendes Muster
        property_type: "safety", "liveness", oder "reachability"
    
    Returns:
        True falls Eigenschaft erfuellt
    """
    pattern_sig = compute_pattern_signature(pattern)
    
    if property_type == "safety":
        # AG(not pattern)
        for state in model.states:
            state_sig = compute_state_signature(state)
            if subgraph_match(pattern_sig, state_sig):
                return False  # Safety verletzt
        return True
    
    elif property_type == "reachability":
        # EF(pattern)
        visited = set()
        queue = [model.initial_state]
        
        while queue:
            state = queue.pop(0)
            if state in visited:
                continue
            visited.add(state)
            
            state_sig = compute_state_signature(state)
            if subgraph_match(pattern_sig, state_sig):
                return True  # Pattern gefunden
            
            queue.extend(model.successors(state))
        
        return False
    
    elif property_type == "liveness":
        # AG(EF(pattern))
        for state in model.states:
            if not check_ef_from_state(model, state, pattern_sig):
                return False
        return True
\end{lstlisting}

\subsection{Zusammenfassung und Ausblick}

Der Subgraph Algorithmus bietet erhebliche Laufzeitverbesserungen
für das temporale Model Checking von PDDL-Planungsproblemen:

\textbf{Hauptergebnisse:}
\begin{enumerate}
	\item \textbf{Polynomiale Komplexität}: Reduktion von $O(n!)$ auf $O(n^3)$
	\item \textbf{Praktische Effizienz}: 15-1000× Speedup gegenüber klassischen Ansätzen
	\item \textbf{Skalierbarkeit}: Anwendbar auf Probleme mit 25+ Objekten
	\item \textbf{Integration}: Nahtlose Kombination mit LTL/CTL Model Checking
\end{enumerate}

\textbf{Anwendungsbereiche:}
\begin{itemize}
	\item[-] PDDL Plan-Validierung und -Optimierung
	\item[-] Robotik: Aktionsplanung mit Safety-Garantien
	\item[-] Logistik: Optimale Routen mit Constraints
	\item[-] Automatisierte Verifikation von Planungssystemen
\end{itemize}

\textbf{Zukünftige Erweiterungen:}
\begin{enumerate}
	\item \textbf{Symbolische Darstellung}: Kombination mit BDDs für noch größere Zustandsräume
	\item \textbf{Parallele Implementierung}: GPU-Beschleunigung der Signatur-Berechnung
	\item \textbf{Erweiterte Patterns}: Unterstützung für zeitliche und probabilistische Constraints
	\item \textbf{Tool-Integration}: Anbindung an Fast Downward, PDDL4J und andere Planner
	\item \textbf{Automatisches Pattern-Mining}: Lernen von Patterns aus erfolgreichen Plänen
\end{enumerate}

Der Subgraph Algorithmus stellt damit einen wichtigen Beitrag zur effizienten Verifikation von Planungssystemen dar und eröffnet neue Möglichkeiten für die praktische Anwendung formaler Methoden in der automatisierten Planung.

\section{Erweiterte PDDL-Beispiele}

Nachdem im vorherigen Kapitel der Subgraph Algorithmus als effiziente Methode für das Model Checking von PDDL-Planungsproblemen vorgestellt wurde, präsentiert dieses Kapitel zwei weitere klassische Planungsdomänen mit vollständigen Beispielen. Diese demonstrieren die Vielseitigkeit von PDDL und bieten zusätzliche Anwendungsfälle für die in Kapitel~7 vorgestellten Verifikationstechniken. Alle Beispiele können mit dem Subgraph Algorithmus effizient verifiziert werden, wobei die strukturellen Muster der jeweiligen Zustandsräume optimal ausgenutzt werden.

\subsection{Logistics-Domäne}

Die Logistics-Domäne modelliert Transport von Paketen zwischen Städten.

\subsubsection{Domänendefinition}

\begin{lstlisting}[language=PDDL, caption=Logistics Domäne]
(define (domain logistics)
  (:requirements :strips :typing)
  
  (:types
    location locatable - object
    package vehicle - locatable
    truck airplane - vehicle
    city airport - location)
  
  (:predicates
    (at ?obj - locatable ?loc - location)
    (in ?pkg - package ?veh - vehicle)
    (in-city ?loc - location ?city - city))
  
  (:action load-truck
    :parameters (?pkg - package ?truck - truck ?loc - location)
    :precondition (and
      (at ?truck ?loc)
      (at ?pkg ?loc))
    :effect (and
      (not (at ?pkg ?loc))
      (in ?pkg ?truck)))
  
  (:action load-airplane
    :parameters (?pkg - package ?airplane - airplane ?loc - airport)
    :precondition (and
      (at ?pkg ?loc)
      (at ?airplane ?loc))
    :effect (and
      (not (at ?pkg ?loc))
      (in ?pkg ?airplane)))
  
  (:action unload-truck
    :parameters (?pkg - package ?truck - truck ?loc - location)
    :precondition (in ?pkg ?truck)
    :effect (and
      (not (in ?pkg ?truck))
      (at ?pkg ?loc)
      (at ?truck ?loc)))
  
  (:action unload-airplane
    :parameters (?pkg - package ?airplane - airplane ?loc - airport)
    :precondition (in ?pkg ?airplane)
    :effect (and
      (not (in ?pkg ?airplane))
      (at ?pkg ?loc)
      (at ?airplane ?loc)))
  
  (:action drive-truck
    :parameters (?truck - truck ?from - location ?to - location ?city - city)
    :precondition (and
      (at ?truck ?from)
      (in-city ?from ?city)
      (in-city ?to ?city))
    :effect (and
      (not (at ?truck ?from))
      (at ?truck ?to)))
  
  (:action fly-airplane
    :parameters (?airplane - airplane ?from - airport ?to - airport)
    :precondition (at ?airplane ?from)
    :effect (and
      (not (at ?airplane ?from))
      (at ?airplane ?to))))
\end{lstlisting}

\subsubsection{Probleminstanz}

\begin{lstlisting}[language=PDDL, caption=Logistics Problem]
(define (problem logistics-problem-1)
  (:domain logistics)
  
  (:objects
    package1 package2 - package
    truck1 truck2 - truck
    airplane1 - airplane
    berlin-tegel hamburg-airport - airport
    berlin-center hamburg-harbor - location
    berlin hamburg - city)
  
  (:init
    ; Package-Positionen
    (at package1 berlin-center)
    (at package2 hamburg-harbor)
    
    ; Fahrzeug-Positionen
    (at truck1 berlin-center)
    (at truck2 hamburg-harbor)
    (at airplane1 berlin-tegel)
    
    ; Stadt-Zuordnungen
    (in-city berlin-center berlin)
    (in-city berlin-tegel berlin)
    (in-city hamburg-harbor hamburg)
    (in-city hamburg-airport hamburg))
  
  (:goal (and
    (at package1 hamburg-harbor)
    (at package2 berlin-center))))
\end{lstlisting}

\subsubsection{Beispielplan}

Ein möglicher Plan zum Transport der Pakete:

\begin{enumerate}
	\item \texttt{load-truck(package1, truck1, berlin-center)}
	\item \texttt{drive-truck(truck1, berlin-center, berlin-tegel, berlin)}
	\item \texttt{unload-truck(package1, truck1, berlin-tegel)}
	\item \texttt{load-airplane(package1, airplane1, berlin-tegel)}
	\item \texttt{fly-airplane(airplane1, berlin-tegel, hamburg-airport)}
	\item \texttt{unload-airplane(package1, airplane1, hamburg-airport)}
	\item \texttt{load-truck(package1, truck2, hamburg-airport)}
	\item \texttt{drive-truck(truck2, hamburg-airport, hamburg-harbor, hamburg)}
	\item \texttt{unload-truck(package1, truck2, hamburg-harbor)}
	\item \texttt{load-truck(package2, truck2, hamburg-harbor)}
	\item \texttt{drive-truck(truck2, hamburg-harbor, hamburg-airport, hamburg)}
	\item \texttt{unload-truck(package2, truck2, hamburg-airport)}
	\item \texttt{load-airplane(package2, airplane1, hamburg-airport)}
	\item \texttt{fly-airplane(airplane1, hamburg-airport, berlin-tegel)}
	\item \texttt{unload-airplane(package2, airplane1, berlin-tegel)}
	\item \texttt{load-truck(package2, truck1, berlin-tegel)}
	\item \texttt{drive-truck(truck1, berlin-tegel, berlin-center, berlin)}
	\item \texttt{unload-truck(package2, truck1, berlin-center)}
\end{enumerate}

\subsubsection{Temporale Eigenschaften}

\begin{beispiel}[Logistics Temporale Eigenschaften]
	\begin{enumerate}
		\item Pakete werden zugestellt
		($\mathit{pkg}_1$ an $\texttt{hamburg-harbor}$, abgekürzt $\mathit{hh}$):
		\begin{align}
		\mathbf{AG}\bigl(\mathrm{at}(\mathit{pkg}_1,\,\mathit{hh})\bigr)
		\end{align}
		
		\item Flugzeuge sind immer an Flughäfen:
		\begin{align}
		\mathbf{AG}(\exists a: \text{airport}: \text{at}(\text{airplane1}, a))
		\end{align}
		
		\item Ein Paket ist nie gleichzeitig an zwei Orten:
		\begin{align}
		\mathbf{AG}\neg(\text{at}(\text{package1}, l_1) \land \text{at}(\text{package1}, l_2) \land l_1 \neq l_2)
		\end{align}
	\end{enumerate}
\end{beispiel}

\subsection{Gripper-Domäne}

Die Gripper-Domäne modelliert einen Roboter mit zwei Greifern, der Bälle zwischen Räumen transportiert.

\subsubsection{Domänendefinition}

\begin{lstlisting}[language=PDDL, caption=Gripper Domäne]
(define (domain gripper)
  (:requirements :strips :typing)
  
  (:types room ball gripper)
  
  (:predicates
    (at-robby ?r - room)
    (at-ball ?b - ball ?r - room)
    (free ?g - gripper)
    (carry ?b - ball ?g - gripper))
  
  (:action move
    :parameters (?from - room ?to - room)
    :precondition (at-robby ?from)
    :effect (and
      (at-robby ?to)
      (not (at-robby ?from))))
  
  (:action pick
    :parameters (?b - ball ?r - room ?g - gripper)
    :precondition (and
      (at-ball ?b ?r)
      (at-robby ?r)
      (free ?g))
    :effect (and
      (carry ?b ?g)
      (not (at-ball ?b ?r))
      (not (free ?g))))
  
  (:action drop
    :parameters (?b - ball ?r - room ?g - gripper)
    :precondition (and
      (carry ?b ?g)
      (at-robby ?r))
    :effect (and
      (at-ball ?b ?r)
      (free ?g)
      (not (carry ?b ?g)))))
\end{lstlisting}

\subsubsection{Probleminstanz}

\begin{lstlisting}[language=PDDL, caption=Gripper Problem]
(define (problem gripper-problem-1)
  (:domain gripper)
  
  (:objects
    room1 room2 - room
    ball1 ball2 ball3 - ball
    left right - gripper)
  
  (:init
    (at-robby room1)
    (free left)
    (free right)
    (at-ball ball1 room1)
    (at-ball ball2 room1)
    (at-ball ball3 room1))
  
  (:goal (and
    (at-ball ball1 room2)
    (at-ball ball2 room2)
    (at-ball ball3 room2))))
\end{lstlisting}

\subsubsection{Optimaler Plan}

Ein optimaler Plan mit 11 Aktionen:

\begin{enumerate}
	\item \texttt{pick(ball1, room1, left)}
	\item \texttt{pick(ball2, room1, right)}
	\item \texttt{move(room1, room2)}
	\item \texttt{drop(ball1, room2, left)}
	\item \texttt{drop(ball2, room2, right)}
	\item \texttt{move(room2, room1)}
	\item \texttt{pick(ball3, room1, left)}
	\item \texttt{move(room1, room2)}
	\item \texttt{drop(ball3, room2, left)}
\end{enumerate}

Dies nutzt optimal beide Greifer, um zwei Bälle gleichzeitig zu transportieren.

\subsubsection{Temporale Analyse}

\begin{beispiel}[Gripper Temporale Eigenschaften]
	\begin{enumerate}
		\item Der Roboter kann beide Greifer gleichzeitig nutzen:
		\begin{align}
		\mathbf{EF}(\exists b_1, b_2, g_1, g_2: \text{carry}(b_1, g_1) \land \text{carry}(b_2, g_2) \land g_1 \neq g_2)
		\end{align}
		
		\item Alle Bälle werden schließlich transportiert:
		\begin{align}
		\mathbf{AF}(\text{at-ball}(\text{ball1}, \text{room2}) \land \text{at-ball}(\text{ball2}, \text{room2}) \land \text{at-ball}(\text{ball3}, \text{room2}))
		\end{align}
		
		\item Ein Greifer trägt maximal einen Ball:
		\begin{align}
		\mathbf{AG}\neg(\text{carry}(b_1, g) \land \text{carry}(b_2, g) \land b_1 \neq b_2)
		\end{align}
	\end{enumerate}
\end{beispiel}

\section{PDDL Planner und Algorithmen}

Dieses Kapitel gibt einen Überblick über gängige PDDL-Planungsalgorithmen. Die hier beschriebenen Verfahren sind komplementär zum Subgraph Algorithmus aus Kapitel~7: Planer generieren Lösungen, während der Subgraph Algorithmus deren Korrektheit und Optimalität effizient verifiziert. Diese Synergie ermöglicht einen vollständigen Workflow von der Plangenerierung über die Verifikation bis zur Zertifizierung.

\subsection{Klassische Planungsalgorithmen}

\subsubsection{Forward Search (State-Space Search)}

Forward Search expandiert vom Initialzustand ausgehend:

\begin{algorithm}[H]
	\caption{Vorwärts-Zustandsraumsuche}
	\begin{algorithmic}[1]
		\Prozedur{ForwardSearch}{$I, G, A$}
		\State $\text{Open} \gets \{I\}$ \Comment{Initialzustand}
		\State $\text{Closed} \gets \emptyset$
		\Waehrend{$\text{Open} \neq \emptyset$}
		\State $s \gets \text{Pop}(\text{Open})$ \Comment{Wähle Zustand}
		\Wenn{$s \models G$}
		\State \Return \textsc{ExtractPlan}($s$)
		\EndeWenn
		\State $\text{Closed} \gets \text{Closed} \cup \{s\}$
		\Fuer{jede Aktion $a \in A$ anwendbar in $s$}
		\State $s' \gets \text{Apply}(a, s)$
		\Wenn{$s' \notin \text{Closed}$}
		\State $\text{Open} \gets \text{Open} \cup \{s'\}$
		\EndeWenn
		\EndeFuer
		\EndeWaehrend
		\State \Return \textbf{Fehler}
		\EndProzedur
	\end{algorithmic}
\end{algorithm}

\subsubsection{Backward Search (Regression)}

Backward Search startet vom Ziel und arbeitet rückwärts:

\begin{algorithm}[H]
	\caption{Rückwärts-Regressionssuche}
	\begin{algorithmic}[1]
		\Prozedur{BackwardSearch}{$I, G, A$}
		\State $\text{Open} \gets \{G\}$ \Comment{Zielzustand}
		\State $\text{Closed} \gets \emptyset$
		\Waehrend{$\text{Open} \neq \emptyset$}
		\State $g \gets \text{Pop}(\text{Open})$
		\Wenn{$I \models g$}
		\State \Return \textsc{ExtractPlan}($g$)
		\EndeWenn
		\State $\text{Closed} \gets \text{Closed} \cup \{g\}$
		\Fuer{jede Aktion $a \in A$}
		\State $g' \gets \text{Regress}(a, g)$ \Comment{Rückwärts-Anwendung}
		\Wenn{$g' \neq \bot$ und $g' \notin \text{Closed}$}
		\State $\text{Open} \gets \text{Open} \cup \{g'\}$
		\EndeWenn
		\EndeFuer
		\EndeWaehrend
		\State \Return \textbf{Fehler}
		\EndProzedur
	\end{algorithmic}
\end{algorithm}

\subsection{Heuristische Planung}

\subsubsection{GraphPlan}

GraphPlan konstruiert einen Planning Graph und extrahiert Pläne:

\begin{enumerate}
	\item Konstruiere Planning Graph Level für Level
	\item Prüfe, ob Ziel erreichbar ist (keine Mutex-Konflikte)
	\item Extrahiere Plan durch Rückwärtssuche
\end{enumerate}

Der Planning Graph alterniert zwischen Proposition-Leveln und Action-Leveln:
\begin{align}
P_0 \xrightarrow{A_0} P_1 \xrightarrow{A_1} P_2 \xrightarrow{A_2} \cdots
\end{align}

\subsubsection{Fast-Forward (FF) Planner}

FF verwendet die relaxed Planning Graph Heuristik:

\begin{definition}[Relaxed Planning Graph]
	Ein relaxed Planning Graph ignoriert Delete-Effekte von Aktionen. Die Heuristik $h^{FF}(s)$ ist die Länge des Lösungsplans im relaxierten Problem.
\end{definition}

Der FF-Planner kombiniert:
\begin{itemize}
	\item[-] Enforced Hill-Climbing: Greedy-Suche mit Heuristik
	\item[-] Best-First Search: Fallback bei Plateaus
\end{itemize}

\subsection{SAT-basierte Planung}

SAT-basierte Planung kodiert Planungsprobleme als aussagenlogische Formeln.

\subsubsection{Bounded Planning}

Für Planlänge $n$ kodiere:
\begin{enumerate}
	\item Initialzustand: $I$
	\item Aktionen und Transitionen: $\bigwedge_{i=0}^{n-1} T_i$
	\item Zielzustand: $G_n$
\end{enumerate}

Die Gesamt-Formel ist:
\begin{align}
\Phi_n = I \land \left(\bigwedge_{i=0}^{n-1} T_i\right) \land G_n
\end{align}

Wenn $\Phi_n$ erfüllbar ist, existiert ein Plan der Länge $n$.

\subsubsection{Inkrementelle SAT-Planung}

\begin{algorithm}[H]
	\caption{SAT-basierte Planung}
	\begin{algorithmic}[1]
		\Prozedur{SATPlan}{$I, G, A$}
		\State $n \gets 0$
		\Wiederhole
		\State $\Phi_n \gets \text{Encode}(I, A, G, n)$
		\State $\text{result} \gets \text{SAT-Solver}(\Phi_n)$
		\Wenn{$\text{result} = \text{SAT}$}
		\State \Return \textsc{ExtractPlan}($\text{result}$)
		\EndeWenn
		\State $n \gets n + 1$
		\Bis{$n > \text{MaxLength}$}
		\State \Return \textbf{Fehler}
		\EndProzedur
	\end{algorithmic}
\end{algorithm}

\subsection{Vergleich der Algorithmen}

\begin{table}[h]
	\centering
	\caption{Vergleich von Planungsalgorithmen}
	\label{tab:planner_comparison}
	\begin{tabular}{llll}
		\toprule
		Algorithmus & Typ & Optimalität & Komplexität \\
		\midrule
		Forward Search & State-Space & Nein & Exponentiell \\
		Backward Search & Regression & Nein & Exponentiell \\
		GraphPlan & Graph-basiert & Ja & Polynomial (pro Level) \\
		FF Planner & Heuristisch & Nein & Praktisch effizient \\
		SAT Planning & Kompilierung & Ja & NP-vollständig \\
		\bottomrule
	\end{tabular}
\end{table}

\subsection{Integration mit Subgraph-basierter Verifikation}

Die vorgestellten Planungsalgorithmen können optimal mit dem Subgraph Algorithmus aus Kapitel~7 kombiniert werden. Die Verifikation mit dem Subgraph Algorithmus fügt dabei nur polynomialen Overhead hinzu, der in der Praxis vernachlässigbar ist:

\begin{table}[h]
	\centering
	\caption{Planungsalgorithmen kombiniert mit Subgraph-Verifikation}
	\label{tab:planner_with_verification}
	\begin{tabular}{llll}
		\toprule
		Algorithmus & Planung & Verifikation & Gesamt \\
		\midrule
		Forward + Subgraph & Exponentiell & $O(|\pi| \cdot n^3)$ & Praktikabel \\
		FF + Subgraph & Heuristisch & $O(|\pi| \cdot n^3)$ & Sehr effizient \\
		SAT + Subgraph & NP-vollständig & $O(|\pi| \cdot n^3)$ & Optimal verifiziert \\
		\bottomrule
	\end{tabular}
\end{table}

Der empfohlene Workflow umfasst fünf Schritte: (1.)~Plangenerierung mit dem FF-Planner oder einem SAT-basierten Ansatz, (2.)~Signaturextraktion aus der Zustandssequenz in $O(n^2)$, (3.)~Pattern-Verifikation für Safety- und Liveness-Eigenschaften, (4.)~Optimalitätsprüfung durch Vergleich der Planlänge mit dem theoretischen Minimum sowie (5.)~Ausgabe eines formal zertifizierten Plans.

\section{Model Checking von PDDL-Domänen}

Dieses Kapitel integriert die in den vorherigen Kapiteln entwickelten Konzepte zu einem vollständigen Verifikations-Workflow. Insbesondere wird gezeigt, wie der Subgraph Algorithmus aus Kapitel~7 in praktische Tools und Workflows eingebettet werden kann, um PDDL-Domänen automatisiert zu verifizieren.

\subsection{Automatische Verifikation mit Subgraph-Integration}

\subsubsection{Erweiterter Workflow}

\begin{enumerate}
	\item \textbf{PDDL zu Kripke}: Konvertiere PDDL-Problem in Kripke-Struktur
	\item \textbf{Signaturberechnung}: Berechne Zustandssignaturen
	      mit $O(n^2)$ pro Zustand (Subgraph Algorithmus)
	\item \textbf{Spezifikation}: Formuliere Eigenschaften als Patterns und temporale Formeln
	\item \textbf{Pattern-Matching}: Nutze Subgraph Algorithmus für strukturelle Eigenschaften in $O(|S| \cdot n^3)$
	\item \textbf{CTL/LTL Model Checking}: Nutze klassische Algorithmen für komplexere temporale Formeln
	\item \textbf{Gegenbeispiel}: Bei Verletzung, analysiere Trace
\end{enumerate}

\begin{table}[h]
	\centering
	\caption{Verifikationszeit mit und ohne Subgraph Algorithmus}
	\label{tab:verification_speedup}
	\begin{tabular}{lrrr}
		\toprule
		Domäne & Klassisch & Mit Subgraph & Speedup \\
		\midrule
		Blocksworld-5 & 2.3s & 0.08s & 29$\times$ \\
		Logistics-4 & 5.1s & 0.15s & 34$\times$ \\
		Gripper-6 & 3.8s & 0.11s & 35$\times$ \\
		\bottomrule
	\end{tabular}
\end{table}

\subsubsection{Workflow-Beispiel für Blocksworld mit Subgraph Algorithmus}

\begin{beispiel}[Blocksworld Verifikation mit Subgraph Algorithmus]
	\textbf{Schritt 1}: Erzeuge Zustandsraum aus PDDL (13 erreichbare Zustände für 3 Blöcke)

	\textbf{Schritt 2}: Berechne Signaturen für alle 13 Zustände: $13 \times O(n^2) = 13 \times O(9) < 0{,}001$s

	\textbf{Schritt 3}: Spezifiziere Eigenschaften:
	\begin{itemize}
		\item[-] Safety-Pattern: $\mathbf{AG}\neg\text{BadState}$
		\item[-] Reachability: $\mathbf{EF}\text{Goal}$
		\item[-] Liveness: $\mathbf{AG}(\text{Request} \to \mathbf{AF}\text{Response})$
	\end{itemize}

	\textbf{Schritt 4}: Pattern-Matching mit Subgraph Algorithmus: $13 \times O(n^3) = O(351)$ Operationen, Zeit $< 0{,}002$s

	\textbf{Schritt 5}: Bei Verletzung, extrahiere Gegenbeispiel-Pfad und konvertiere zurück zu PDDL-Aktionssequenz

	\textbf{Gesamtzeit}: 0{,}006s (vs. 0{,}18s klassisch) $= 30\times$ Speedup. Alle Eigenschaften verifiziert: $\checkmark$
\end{beispiel}

\subsection{Praktische Tools}

\subsubsection{PDDL4J}

PDDL4J ist eine Java-Bibliothek für PDDL:
\begin{itemize}
	\item[-] Parser für PDDL 3.0
	\item[-] Mehrere Planner (FF, HSP, etc.)
	\item[-] API für eigene Algorithmen
\end{itemize}

\subsubsection{Fast Downward}

Fast Downward ist ein state-of-the-art Planner:
\begin{itemize}
	\item[-] Unterstützt PDDL 2.2
	\item[-] Verschiedene Heuristiken
	\item[-] Exzellente Performance
\end{itemize}

\subsubsection{Integration mit NuSMV}

NuSMV ist ein symbolischer Model Checker:

\begin{enumerate}
	\item Konvertiere PDDL zu SMV-Format
	\item Definiere CTL-Eigenschaften in SMV
	\item Führe NuSMV aus
	\item Analysiere Ergebnisse
\end{enumerate}

\subsection{Case Study: Blocksworld Vollständige Verifikation}

\subsubsection{Zu verifizierende Eigenschaften}

\begin{enumerate}
	\item \textbf{Korrektheit}: Plan erreicht Ziel
	\begin{align}
	\mathbf{AF}(\text{on}(C, B) \land \text{on}(B, A))
	\end{align}
	
	\item \textbf{Safety}: Niemals inkonsistenter Zustand
	\begin{align}
	\mathbf{AG}(\text{handempty} \lor \exists! x: \text{holding}(x))
	\end{align}
	
	\item \textbf{Effizienz}: Minimale Anzahl Aktionen
	\begin{align}
	\text{optimal}(\pi) \iff |\pi| = \min\{|\pi'| \mid \pi' \text{ ist Plan}\}
	\end{align}
	
	\item \textbf{Determinismus}: Eindeutiger Nachfolgezustand
	\begin{align}
	\mathbf{AG}\mathbf{AX}\varphi \to \mathbf{AG}\mathbf{EX}\varphi
	\end{align}
\end{enumerate}

\subsubsection{Verifikationsergebnisse}

Für das Blocksworld-Problem mit 3 Blöcken:

\begin{table}[h]
	\centering
	\caption{Verifikationsergebnisse Blocksworld}
	\label{tab:blocksworld_verification}
	\begin{tabular}{lll}
		\toprule
		Eigenschaft & Ergebnis & Zeit \\
		\midrule
		Korrektheit & $\checkmark$ Erfüllt & 0.002s \\
		Safety (Hand) & $\checkmark$ Erfüllt & 0.001s \\
		Safety (Blocks) & $\checkmark$ Erfüllt & 0.001s \\
		Optimalität & $\checkmark$ Optimal (4 Schritte) & 0.003s \\
		Determinismus & $\checkmark$ Erfüllt & 0.001s \\
		\bottomrule
	\end{tabular}
\end{table}

\section{Hybride Verifikation: Kontinuierliche Systeme und diskrete Abstraktion}

Die bisher entwickelten Verifikationsverfahren operieren ausschließlich auf diskreten Zustandsräumen. Reale Systeme hingegen -- Roboter, autonome Fahrzeuge, Produktionsanlagen, medizinische Geräte -- sind cyber-physische Systeme, deren Zustand durch kontinuierliche Größen wie Position, Geschwindigkeit, Temperatur oder Druck beschrieben wird. Die Verbindung zwischen der kontinuierlichen physikalischen Welt und der diskreten Verifikationswelt erfordert eine sorgfältige Abstraktion, die unvermeidlich Fehler einführt. Dieses Kapitel untersucht systematisch, wie diese Abstraktionsfehler entstehen, wie sie formal charakterisiert werden können, wie sie durch Kalibrierung minimiert werden und ob es Klassen realer Systeme gibt, für die eine vollständige Fehlervermeidung beweisbar erreichbar ist.

\subsection{Hybride Automaten als formales Modell}

\begin{definition}[Hybrider Automat]
Ein hybrider Automat ist ein Tupel $H$:
\begin{align*}
  H = (Q,\; X,\; f,\; \mathit{Init},\; \mathit{Inv},\; E,\; G,\; R)
\end{align*}
mit folgenden Komponenten:
\begin{itemize}
  \item[-] $Q$: endliche Menge diskreter Zustände (Moden)
  \item[-] $X \subseteq \mathbb{R}^n$: kontinuierlicher Zustandsraum
  \item[-] $f: Q \times X \to \mathbb{R}^n$: Vektorfeld, beschreibt die kontinuierliche Dynamik $\dot{x} = f(q, x)$ in Modus $q$
  \item[-] $\mathit{Init} \subseteq Q \times X$: Initialbedingung
  \item[-] $\mathit{Inv}: Q \to 2^X$: Invariante, legt fest, in welchen kontinuierlichen Zuständen Modus $q$ zulässig ist
  \item[-] $E \subseteq Q \times Q$: diskrete Übergänge (Kanten)
  \item[-] $G: E \to 2^X$: Schaltbedingung (Guard) für jeden Übergang
  \item[-] $R: E \times X \to 2^X$: Rücksetzabbildung (Reset) bei Übergängen
\end{itemize}
\end{definition}

Ein hybrider Automat kombiniert das Schaltverhalten eines endlichen Automaten mit der kontinuierlichen Dynamik eines Differentialgleichungssystems. Reale cyber-physische Systeme lassen sich natürlich in diesem Formalismus modellieren: Ein Roboter befindet sich in verschiedenen Betriebsmodi (Fahren, Bremsen, Stehen), und innerhalb jedes Modus beschreibt ein Differentialgleichungssystem die Bewegungsdynamik.

\begin{beispiel}[Hybrider Automat: Fahrender Roboter]
Ein mobiler Roboter besitzt drei Moden: $Q = \{\mathit{Fahren}, \mathit{Bremsen}, \mathit{Stehen}\}$. Der kontinuierliche Zustand ist $(x, v) \in \mathbb{R}^2$ mit Position $x$ und Geschwindigkeit $v$. Die Vektorfelder sind:
\begin{align}
  f(\mathit{Fahren}, (x,v))  &= (v,\; a_{\max}) \\
  f(\mathit{Bremsen}, (x,v)) &= (v,\; -a_{\max}) \\
  f(\mathit{Stehen}, (x,v))  &= (0,\; 0)
\end{align}
Der Übergang von \textit{Fahren} nach \textit{Bremsen} wird ausgelöst, wenn ein Hindernis in Distanz $d_{\mathit{safe}}$ erkannt wird (Guard: $x_{\mathit{obstacle}} - x \leq d_{\mathit{safe}}$); das Reset setzt $v$ nicht zurück, da die Bremsung kontinuierlich beginnt.
\end{beispiel}

\subsection{Diskretisierung und Abstraktionsfehler}

Die Umwandlung eines hybriden Automaten $H$ in eine endliche Kripke-Struktur $M$ erfolgt durch \emph{Diskretisierung}: Der kontinuierliche Zustandsraum $X$ wird in endlich viele Zellen $C_1, \ldots, C_k$ eingeteilt, und jede Zelle repräsentiert einen diskreten Zustand.

\begin{definition}[Gitterbasierte Diskretisierung]
Für eine Auflösung $\delta > 0$ und einen kompakten Zustandsraum $X = [a_1, b_1] \times \cdots \times [a_n, b_n]$ definieren wir die Gitterzellen:
\begin{align}
  C_{i_1, \ldots, i_n} = \bigl[a_1 + i_1 \delta,\; a_1 + (i_1+1)\delta\bigr) \times \cdots \times \bigl[a_n + i_n \delta,\; a_n + (i_n+1)\delta\bigr)
\end{align}
für $i_j \in \{0, 1, \ldots, \lceil(b_j - a_j)/\delta\rceil - 1\}$. Die Anzahl der Gitterzellen ist $N(\delta) = \prod_{j=1}^n \lceil(b_j - a_j)/\delta\rceil$.
\end{definition}

Diese Diskretisierung führt zu zwei grundlegenden Fehlerquellen:

\textbf{Räumlicher Abstraktionsfehler.} Alle Punkte einer Zelle werden mit dem Zellmittelpunkt~$\hat{x}_{C}$ identifiziert. Der maximale räumliche Fehler pro Zelle beträgt:
\begin{align}
  \epsilon_{\mathit{räumlich}}(\delta) = \frac{\sqrt{n} \cdot \delta}{2}
\end{align}
wobei $n$ die Dimension des Zustandsraums ist ($\ell_2$-Norm zum Zellmittelpunkt).

\textbf{Dynamischer Propagationsfehler.} Wird die Dynamik $\dot{x} = f(q, x)$ durch einen diskreten Zeitschritt $\Delta t$ approximiert (Euler-Methode: $x_{k+1} = x_k + \Delta t \cdot f(q, x_k)$), so akkumuliert sich ein Fehler über die Zeit. Für Lipschitz-stetige Vektorfelder gilt:

\begin{theorem}[Fehlerpropagation bei Euler-Diskretisierung]
Sei $f(q,\cdot)$ Lipschitz-stetig mit Konstante $L>0$. Nach $k$ Euler-Schritten der Weite~$\Delta t$ gilt für den Fehler zur exakten Lösung~$x^*(t)$:
\begin{align}
  \|x^*(k\Delta t) - \hat{x}_k\| \leq \frac{M \Delta t}{2L} \bigl(e^{Lk\Delta t} - 1\bigr)
\end{align}
wobei $M = \sup_{q,x} \|\ddot{x}^*(t)\|$ die maximale zweite Ableitung der exakten Lösung beschränkt.
\end{theorem}

\begin{proof}
Der lokale Abbruchfehler des Euler-Verfahrens pro Schritt beträgt durch Taylor-Entwicklung:
\begin{align}
  \|x^*(t + \Delta t) - (x^*(t) + \Delta t \cdot f(q, x^*(t)))\| \leq \frac{M}{2} (\Delta t)^2.
\end{align}
Sei $e_k = \|x^*(k\Delta t) - \hat{x}_k\|$ der globale Fehler nach $k$ Schritten, mit $e_0 = 0$. Es gilt:
\begin{align}
  e_{k+1} &\leq \|x^*(( k+1)\Delta t) - x^*(k\Delta t) - \Delta t \cdot f(q, x^*(k\Delta t))\| \\
           &\quad + \|\Delta t \cdot f(q, x^*(k\Delta t)) - \Delta t \cdot f(q, \hat{x}_k)\| + e_k \notag \\
           &\leq \frac{M}{2}(\Delta t)^2 + L \Delta t \cdot e_k + e_k \notag \\
           &= (1 + L\Delta t) e_k + \frac{M}{2}(\Delta t)^2. \notag
\end{align}
Diese Rekurrenz hat die Lösung:
\begin{align}
  e_k \leq \frac{M(\Delta t)^2}{2} \cdot \frac{(1+L\Delta t)^k - 1}{L\Delta t} = \frac{M\Delta t}{2L}\bigl((1+L\Delta t)^k - 1\bigr).
\end{align}
Mit $(1 + L\Delta t)^k \leq e^{Lk\Delta t}$ folgt die Behauptung.
\end{proof}

\begin{remark}
Das Ergebnis zeigt zwei kritische Abhängigkeiten: Der Fehler wächst linear in $\Delta t$ (für kleine $\Delta t$) und exponentiell in der Simulationszeit $T = k\Delta t$. Für große Lipschitz-Konstanten $L$ -- das heißt, für schnell veränderliche Systeme -- ist die exponentielle Fehlerverstärkung besonders ausgeprägt.
\end{remark}

\subsection{Formale Fehlerschranken und Kalibrierung}

Für die praktische Verifikation ist es entscheidend, nicht nur zu wissen, dass ein Fehler existiert, sondern ihn quantitativ zu begrenzen. Wir entwickeln eine systematische Theorie der Kalibrierung, die den Abstraktionsfehler unter eine vorgegebene Schranke $\varepsilon > 0$ drückt.

\subsubsection{Kombinierter Gesamtfehler}

Der Gesamtfehler der Diskretisierung setzt sich aus räumlichem Fehler und Propagationsfehler zusammen:

\begin{theorem}[Gesamtfehler der Diskretisierung]
\label{thm:gesamtfehler}
Für einen hybriden Automaten mit $n$-dimensionalem Zustandsraum, Lipschitz-Konstante $L$, maximaler zweiter Ableitung $M$, Diskretisierungsauflösung $\delta$ und Zeitschrittweite $\Delta t$ gilt für den Gesamtfehler nach Simulationszeit $T$:
\begin{align}
  \epsilon_{\mathit{gesamt}}(\delta, \Delta t, T) \leq \underbrace{\frac{\sqrt{n}\,\delta}{2}}_{\text{räumlich}} + \underbrace{\frac{M\Delta t}{2L}\bigl(e^{LT} - 1\bigr)}_{\text{dynamisch}}
\end{align}
\end{theorem}

\begin{proof}
Der räumliche Fehler beträgt höchstens den halben Zelldurchmesser~$\frac{\sqrt{n}\delta}{2}$ (Mittelpunktsabstand der $n$-di\-men\-sio\-na\-len Würfelzelle der Seitenlänge~$\delta$). Der dynamische Fehler folgt aus Theorem~11.1 mit $k = T/\Delta t$. Da beide Fehlerquellen unabhängig sind, addieren sie sich im schlechtesten Fall.
\end{proof}

\subsubsection{Kalibrierungsregel}

Aus Theorem~\ref{thm:gesamtfehler} lässt sich direkt eine \emph{Kalibrierungsregel} ableiten: Für eine vorgegebene Fehlertoleranz $\varepsilon$ wähle $\delta$ und $\Delta t$ so, dass:
\begin{align}
  \frac{\sqrt{n}\,\delta}{2} \leq \frac{\varepsilon}{2} \quad \Rightarrow \quad \delta \leq \frac{\varepsilon}{\sqrt{n}}
  \label{eq:kalibrierung_raeumlich}
\end{align}
\begin{align}
  \frac{M\Delta t}{2L}(e^{LT}-1) \leq \frac{\varepsilon}{2} \quad \Rightarrow \quad \Delta t \leq \frac{L\varepsilon}{M(e^{LT}-1)}
  \label{eq:kalibrierung_dynamisch}
\end{align}

\begin{definition}[Kalibriertes Diskretisierungsschema]
Ein Schema $(\delta^*, \Delta t^*)$ heißt \emph{$\varepsilon$-kalibriert} für Parameter $(n, L, M, T)$, wenn gilt:
\begin{align}
  \delta^* = \frac{\varepsilon}{\sqrt{n}}, \quad \Delta t^* = \frac{L\varepsilon}{M(e^{LT}-1)}.
\end{align}
\end{definition}

\begin{satz}[Korrektheit der Kalibrierung]
Ein $\varepsilon$-kalibriertes Schema garantiert $\epsilon_{\mathit{gesamt}} \leq \varepsilon$.
\end{satz}

\begin{proof}
Direkt durch Einsetzen in Theorem~\ref{thm:gesamtfehler}:
\begin{align}
  \epsilon_{\mathit{gesamt}} \leq \frac{\sqrt{n}}{2} \cdot \frac{\varepsilon}{\sqrt{n}} + \frac{M}{2L} \cdot \frac{L\varepsilon}{M(e^{LT}-1)} \cdot (e^{LT}-1) = \frac{\varepsilon}{2} + \frac{\varepsilon}{2} = \varepsilon. \qquad \square
\end{align}
\end{proof}

\subsubsection{Adaptives Kalibrierungsverfahren}

In der Praxis sind die Systemparameter $(L, M)$ oft nicht exakt bekannt und müssen empirisch geschätzt werden. Das folgende adaptive Verfahren passt $\delta$ und $\Delta t$ zur Laufzeit an:

\begin{algorithm}[H]
  \caption{Adaptives Kalibrierungsverfahren}
  \begin{algorithmic}[1]
    \Eingabe Hybrider Automat $H$, Fehlerziel $\varepsilon$, Schätzintervall $T_{\mathit{est}}$
    \Ausgabe Kalibriertes Schema $(\delta^*, \Delta t^*)$
    \Prozedur{AdaptiveKalibrierung}{$H, \varepsilon, T_{\mathit{est}}$}
      \State Simuliere $H$ für Dauer $T_{\mathit{est}}$ mit grober Schrittweite $\Delta t_0$
      \State Schätze $\hat{L} \gets \max_{t} \|J_f(q(t), x(t))\|$ \Comment{Jacobi-Norm als Lipschitz-Schätzung}
      \State Schätze $\hat{M} \gets \max_{t} \|\ddot{x}(t)\|$ \Comment{Numerische zweite Ableitung}
      \State $\delta^* \gets \varepsilon / \sqrt{n}$
      \State $\Delta t^* \gets \hat{L}\varepsilon \,/\, (\hat{M}(e^{\hat{L}T_{\mathit{est}}}-1))$
      \State \Return $(\delta^*, \Delta t^*)$
    \EndProzedur
  \end{algorithmic}
\end{algorithm}

Das adaptive Verfahren hat den Vorteil, dass es keine a-priori-Kenntnis der Systemparameter voraussetzt. Die Schätzung der Lipschitz-Konstante $\hat{L}$ über die Jacobi-Norm ist eine konservative obere Schranke und gewährleistet damit, dass der resultierende Fehler das Ziel $\varepsilon$ nicht überschreitet.

\subsection{Kalibrierung für konkrete Systemklassen}

Das folgende zentrale Ergebnis dieses Kapitels zeigt, dass für bestimmte Klassen realer Systeme nicht nur eine Fehlerminimierung, sondern -- unter präzise formulierten Bedingungen -- eine vollständige Fehlervermeidung im Verifikationskontext erreichbar ist.

\subsubsection{Klasse 1: Stückweise lineare Systeme}

\begin{definition}[Stückweise lineares System]
Ein hybrides System heißt \emph{stückweise linear}, wenn in jedem Modus $q \in Q$ gilt:
\begin{align}
  f(q, x) = A_q x + b_q, \quad A_q \in \mathbb{R}^{n \times n},\; b_q \in \mathbb{R}^n.
\end{align}
\end{definition}

Für stückweise lineare Systeme lässt sich die exakte Lösung der Differentialgleichung in geschlossener Form angeben:
\begin{align}
  x(t) = e^{A_q t} x_0 + A_q^{-1}(e^{A_q t} - I) b_q \quad \text{(für } A_q \text{ invertierbar)}.
\end{align}

\begin{theorem}[Fehlerfreie Diskretisierung stückweise linearer Systeme]
\label{thm:fehlerfrei_linear}
Für ein stückweise lineares System mit bekannten Matrizen $A_q$ existiert ein Diskretisierungsschema, das den dynamischen Propagationsfehler exakt auf null reduziert: Wähle als Übergangsrelation der diskreten Kripke-Struktur die durch die exakte Matrixexponential-Lösung induzierte Abbildung $x \mapsto e^{A_q \Delta t} x + A_q^{-1}(e^{A_q \Delta t} - I) b_q$.
\end{theorem}

\begin{proof}
Die Euler-Approximation wird durch die exakte Lösung ersetzt. Der lokale Abbruchfehler ist dann identisch null, da die Übergangsrelation der diskreten Struktur aus der analytischen Lösung der Differentialgleichung abgeleitet ist. Der räumliche Fehler $\frac{\sqrt{n}\delta}{2}$ bleibt bestehen, kann aber durch $\delta \to 0$ beliebig klein gemacht werden. Im Grenzfall $\delta \to 0$ verschwindet auch dieser, was jedoch einen unendlichen Zustandsraum zur Folge hat. Für jedes endliche $\varepsilon > 0$ ist jedoch das $\varepsilon$-kalibrierte Schema mit $\delta = \varepsilon/\sqrt{n}$ und dem Matrixexponential als Übergangsfunktion fehlerfrei bezüglich des dynamischen Fehleranteils.
\end{proof}

\textbf{Praktische Relevanz.} Stückweise lineare Systeme modellieren eine Vielzahl realer Systeme: lineare Regelkreise, Roboter mit linearer Kinematik, elektrische Netzwerke im stationären Betrieb. Für alle diese Systeme kann der dynamische Abstraktionsfehler durch die Verwendung der exakten Matrixexponential-Lösung auf null gesetzt werden.

\subsubsection{Klasse 2: Systeme mit beschränktem Zustandsraum und Lipschitz-Dynamik}

\begin{theorem}[Gleichmäßige Kalibrierbarkeit]
\label{thm:gleichmaessig_kalibrierbar}
Sei $\mathcal{S}$ eine Klasse hybrider Systeme, für die eine gemeinsame obere Schranke $L_{\max}$ für alle Lipschitz-Konstanten und eine gemeinsame Schranke $M_{\max}$ für alle zweiten Ableitungen existiert. Dann ist für jedes $\varepsilon > 0$ das Schema
\begin{align}
  \delta^* = \frac{\varepsilon}{\sqrt{n}}, \quad \Delta t^* = \frac{L_{\max}\varepsilon}{M_{\max}(e^{L_{\max}T}-1)}
\end{align}
für \emph{alle} Systeme aus $\mathcal{S}$ gleichzeitig $\varepsilon$-kalibriert.
\end{theorem}

\begin{proof}
Da $L_{\max} \geq L$ und $M_{\max} \geq M$ für alle Systeme der Klasse gilt, ist $\Delta t^* \leq \frac{L\varepsilon}{M(e^{LT}-1)}$ für jedes einzelne System (da $\frac{L_{\max}}{M_{\max}} \leq \frac{L_{\max}}{M_{\max}}$ und die Funktion $L \mapsto \frac{L}{e^{LT}-1}$ für $L > 0$ monoton fallend ist). Damit ist Bedingung~(\ref{eq:kalibrierung_dynamisch}) für alle Systeme der Klasse erfüllt. Bedingung~(\ref{eq:kalibrierung_raeumlich}) hängt nicht von den Systemparametern ab und ist trivialerweise für alle Systeme gleich.
\end{proof}

Dieses Ergebnis ist von grundlegender praktischer Bedeutung: Es besagt, dass ein \emph{einziges} Kalibrierungsschema für eine ganze Klasse physikalischer Systeme verwendet werden kann, ohne jedes System individuell zu kalibrieren -- eine wesentliche Vereinfachung für industrielle Zertifizierungsprozesse.

\subsubsection{Klasse 3: Robotersysteme mit bekannter Kinematik}

Robotersysteme bilden eine besonders wichtige Anwendungsklasse. Ihre Kinematik ist durch die Denavit-Hartenberg-Parameter vollständig beschrieben und damit bekannt. Die Jacobi-Matrix $J(q)$ des Roboters (nicht zu verwechseln mit der Jacobi-Matrix des Vektorfelds) liefert direkt eine Schranke für die Lipschitz-Konstante.

\begin{satz}[Lipschitz-Schranke für Roboterkinematik]
Für einen Roboter mit $m$ Gelenken, Jacobi-Matrix $J(q)$ und maximalem Gelenkwinkel $q_{\max}$ gilt:
\begin{align}
  L \leq \|J(q)\|_2 \cdot \|\dot{q}_{\max}\|
\end{align}
wobei $\|\dot{q}_{\max}\|$ die maximale Gelenkwinkelgeschwindigkeit ist.
\end{satz}

\begin{proof}
Die Endeffektorgeschwindigkeit ist $\dot{x} = J(q)\dot{q}$. Die Lipschitz-Konstante des Vektorfelds $f(q_{\mathit{mode}}, x) = J(q)\dot{q}$ bezüglich $x$ ergibt sich aus der Operator-Norm der Jacobi-Matrix, die die maximale Streckung von Tangentenvektoren beschreibt. Da $\|\dot{q}\| \leq \|\dot{q}_{\max}\|$, folgt $L \leq \|J\|_2 \cdot \|\dot{q}_{\max}\|$.
\end{proof}

Für einen konkreten Industrie-Roboterarm (z.\,B.\ KUKA KR 10 R1100) mit bekannten Kinematikparametern und gegebenen Gelenkwinkelgrenzen lässt sich damit $L_{\max}$ analytisch bestimmen -- und das $\varepsilon$-kalibrierte Schema nach Theorem~\ref{thm:gleichmaessig_kalibrierbar} kann vor dem Deployment berechnet werden.

\begin{beispiel}[Kalibrierung eines KUKA-Roboterarms]
Für einen 6-Achs-Roboter mit $L_{\max} = 12{,}5\,\text{rad/s}$,
$M_{\max} = 80\,\text{rad/s}^2$, $T = 5\,\text{s}$, $n=6$
und $\varepsilon = 0{,}01\,\text{rad}$ ergibt sich:
\begin{align}
  \delta^* &= \frac{0{,}01}{\sqrt{6}} \approx 0{,}0041\,\text{rad}, \\
  \Delta t^* &= \frac{12{,}5 \cdot 0{,}01}{80 \cdot (e^{12{,}5 \cdot 5} - 1)} \approx 2{,}9 \cdot 10^{-8}\,\text{s}.
\end{align}
Die sehr kleine Schrittweite~$\Delta t^*$ zeigt, dass die Matrixexponential-Methode (Klasse~1) die praktikablere Wahl ist, sofern die Kinematik lokal linearisiert werden kann.
\end{beispiel}

\subsection{Soundness der Abstraktion: Formale Ãœberabstraktion}

Bisher haben wir Fehlerschranken für \emph{korrekte} Approximationen untersucht. Für die formale Verifikation ist jedoch eine stärkere Eigenschaft wünschenswert: \emph{Soundness}, d.\,h., dass kein Safety-Fehler des diskreten Modells übersehen wird, auch wenn das Modell mehr Verhalten erlaubt als das reale System.

\begin{definition}[Sound Abstraktion]
Eine diskrete Kripke-Struktur $M$ heißt \emph{sound} bezüglich des hybriden Automaten $H$, wenn für jede Trajektorie $\xi$ von $H$ ein entsprechender Pfad $\pi$ in $M$ existiert, der dieselben Labelbedingungen erfüllt:
\begin{align}
  \forall \xi \in \mathit{Traj}(H),\; \exists\, \pi \in \mathit{Paths}(M)\colon L_H(\xi) \subseteq L_M(\pi).
\end{align}
\end{definition}

Eine sound Abstraktion kann niemals Safety-Eigenschaften fälschlicherweise beweisen: Wenn $M \models \mathbf{AG}\neg\mathit{unsafe}$, dann gilt auch $H \models \mathbf{AG}\neg\mathit{unsafe}$.

\begin{theorem}[Ãœberabstraktion durch $\varepsilon$-Erweiterung]
\label{thm:soundness}
Sei $M_\varepsilon$ die diskrete Kripke-Struktur, die durch $\varepsilon$-kalibrierte Diskretisierung von $H$ entsteht, wobei die Transitionsrelation um alle Zellen erweitert wird, die innerhalb des Fehlerradius $\varepsilon$ einer tatsächlichen Folgezelle liegen:
\begin{align}
  R_\varepsilon(C, C') \iff \exists\, x \in C,\; x' \in C'\colon \|x' - \phi_{\Delta t}(x)\| \leq \varepsilon
\end{align}
wobei $\phi_{\Delta t}$ der exakte Fluss des Systems nach Zeit $\Delta t$ ist. Dann ist $M_\varepsilon$ eine sound Abstraktion von $H$.
\end{theorem}

\begin{proof}
Sei $\xi = (q_0, x_0), (q_1, x_1), \ldots$ eine beliebige Trajektorie von $H$. Nach Definition der $\varepsilon$-Erweiterung gilt für jeden Übergang $(x_k, x_{k+1})$ der Trajektorie: Da $x_{k+1} = \phi_{\Delta t}(x_k)$ (exakter Fluss), existiert die Zelle $C_{k+1}$ mit $x_{k+1} \in C_{k+1}$, und nach der $\varepsilon$-Erweiterung gilt $R_\varepsilon(C_k, C_{k+1})$ für die entsprechenden Zellen. Also hat jede Trajektorie von $H$ einen korrespondierenden Pfad in $M_\varepsilon$.
\end{proof}

Die $\varepsilon$-Erweiterung ist die entscheidende technische Idee: Das diskrete Modell erlaubt \emph{mehr} Übergänge als das reale System (Überabstraktion), vergisst aber keinen Übergang des realen Systems (Soundness). Safety-Beweise im diskreten Modell sind damit direkt übertragbar.

\subsection{Vollständige Fehlervermeidung}

Die bisher entwickelten Methoden minimieren den Fehler, eliminieren ihn aber nur im Grenzfall $\delta, \Delta t \to 0$. Wir untersuchen nun, unter welchen Bedingungen eine vollständige Fehlervermeidung im endlichen Modell möglich ist.

\begin{theorem}[Notwendige Bedingung für fehlerfreie Diskretisierung]
\label{thm:notwendig}
Eine vollständig fehlerfreie endliche Diskretisierung (mit $\epsilon_{\mathit{gesamt}} = 0$ und endlichem $N(\delta)$) ist genau dann möglich, wenn das hybride System $H$ in jedem Modus eine \emph{lineare} Dynamik besitzt und die Guard-Mengen und Invarianten durch Hyperebenen beschrieben werden, die mit dem Diskretisierungsgitter kompatibel sind.
\end{theorem}

\begin{proof}
\textbf{Hinreichende Richtung:} Für lineare Dynamik wird die Matrixexponential-Lösung als Übergangsfunktion genutzt (Theorem~\ref{thm:fehlerfrei_linear}), was den dynamischen Fehler eliminiert. Mit dem Gitter kompatible Hyperebenen-Guards erzeugen keinen Schnittfehler. Alle relevanten Grenzpunkte liegen exakt auf dem Gitter, daher ist $\epsilon_{\mathit{gesamt}} = 0$.

\textbf{Notwendige Richtung:} Ist die Dynamik nichtlinear, so kann die Lösung nicht in geschlossener Form dargestellt werden, und jede endliche Approximation akkumuliert einen positiven lokalen Abbruchfehler. Fallen Guard-Grenzen nicht mit dem Gitter zusammen, entstehen Zellen, die teils im Guard, teils außerhalb liegen; die Klassifikation solcher Randzellen erzwingt einen positiven Klassifikationsfehler.
\end{proof}

\begin{remark}
Theorem~\ref{thm:notwendig} klärt, warum vollständige Fehlervermeidung in der Praxis selten erreichbar ist: Die meisten realen Systeme sind nichtlinear. Es identifiziert jedoch präzise die Klasse der Systeme -- lineare Dynamik mit kompatiblen Guards --, für die sie erreichbar ist. Dies umfasst beispielsweise viele Regelkreise in der Luft- und Raumfahrt (linearisierte Flugdynamik), elektrische Antriebe (lineare Motormodelle) und Teilbereiche der Robotik (linearisierte Kinematik in Arbeitspunktnähe).
\end{remark}

\subsection{Anwendung auf konkrete Systeme}

\subsubsection{TurtleBot3: Mobiler Roboter}

Der TurtleBot3 ist ein verbreiteter mobiler Roboter mit Differentialantrieb. Seine kinematische Gleichung lautet:
\begin{align}
  \dot{x} = v\cos\theta, \quad \dot{y} = v\sin\theta, \quad \dot{\theta} = \omega,
\end{align}
wobei $(x, y, \theta)$ Position und Orientierung sind und $(v, \omega)$ die Steuereingaben. Das System ist nichtlinear (wegen $\cos\theta$ und $\sin\theta$). Allerdings ist die Nichtlinearität beschränkt: Für $\theta \in [-\pi, \pi]$ gilt $|\cos\theta| \leq 1$ und $|\sin\theta| \leq 1$, was eine globale Lipschitz-Konstante von $L = v_{\max}$ liefert.

Mit den typischen Parameterwerten $v_{\max} = 0{,}5\,\text{m/s}$, $\omega_{\max} = 2{,}84\,\text{rad/s}$, $n = 3$ ergibt das $\varepsilon$-kalibrierte Schema für $\varepsilon = 0{,}05\,\text{m}$ und $T = 30\,\text{s}$:
\begin{align}
  \delta^* &= \frac{0{,}05}{\sqrt{3}} \approx 0{,}029\,\text{m}, \\
  \Delta t^* &= \frac{0{,}5 \cdot 0{,}05}{M_{\max}(e^{0{,}5 \cdot 30}-1)} \approx \frac{0{,}025}{M_{\max} \cdot e^{15}}.
\end{align}

Für den TurtleBot3 ergibt sich damit ein Gitter der Auflösung ca. 3\,cm, was praktisch handhabbar ist. Die Verifikation der Kollisionsfreiheit $\mathbf{AG}(d_{\mathit{obstacle}} > d_{\mathit{safe}})$ kann mit dem Subgraph Algorithmus in $O(|S| \cdot n^3)$ durchgeführt werden, wobei $|S|$ die Anzahl der Gitterzellen im verifizierten Bereich ist.

\subsubsection{Autonomes Fahrzeug: Einspurmodell}

Das kinematische Einspurmodell (Fahrradmodell) eines autonomen Fahrzeugs lautet:
\begin{align}
  \dot{x} &= v\cos(\theta + \beta), \\
  \dot{y} &= v\sin(\theta + \beta), \\
  \dot{\theta} &= \frac{v\cos\beta}{l_r}\tan\delta_f,
\end{align}
wobei $\beta = \arctan\!\bigl(\frac{l_r}{l_f + l_r}\tan\delta_f\bigr)$ der Schwimmwinkel, $l_f, l_r$ die Abstände von Vorder- und Hinterachse zum Schwerpunkt und $\delta_f$ der Lenkwinkel sind.

Die Nichtlinearität ist stärker als beim TurtleBot3, aber die Lipschitz-Konstante bleibt beschränkt, da alle trigonometrischen Funktionen auf $[-1, 1]$ beschränkt sind. Für $v_{\max} = 30\,\text{km/h} \approx 8{,}3\,\text{m/s}$ und $n = 4$ (Position, Kurs, Geschwindigkeit) ergibt das Kalibrierungsschema für $\varepsilon = 0{,}1\,\text{m}$ und $T = 10\,\text{s}$:
\begin{align}
  \delta^* = \frac{0{,}1}{\sqrt{4}} = 0{,}05\,\text{m}.
\end{align}
Dieser Wert von 5\,cm räumlicher Auflösung ist mit moderner Sensorik (RTK-GNSS, LiDAR) gut erreichbar und erlaubt eine praxistaugliche Verifikation.

\subsubsection{Industrieroboter: KUKA KR 10}

Für den KUKA KR 10 R1100-Roboter mit 6 Freiheitsgraden wurde gezeigt (Beispiel~11.1), dass die Euler-Diskretisierung eine extrem kleine Schrittweite erfordert. Stattdessen bietet sich die Klasse-1-Methode (Matrixexponential) an: Die Kinematik kann in Arbeitspunktnähe mit der Jacobi-Matrix linearisiert werden:
\begin{align}
  \Delta\dot{x}_e \approx J(q_0)\,\Delta\dot{q},
\end{align}
wobei $q_0$ der aktuelle Gelenkwinkel-Arbeitspunkt ist. In dieser linearisierten Darstellung ist der dynamische Fehler exakt null (Theorem~\ref{thm:fehlerfrei_linear}), und nur der räumliche Fehler $\delta^* = \varepsilon/\sqrt{6}$ verbleibt. Die Linearisierung ist für kleine Bewegungen um den Arbeitspunkt herum -- typisch für Montageprozesse -- eine ausgezeichnete Näherung.

\subsubsection{Drohnen-Schwarm: Universelle Kalibrierung}

Für einen Drohnen-Schwarm mit $k$ identischen Quadrocoptern stellt sich die Frage, ob ein gemeinsames Kalibrierungsschema für alle Drohnen existiert. Nach Theorem~\ref{thm:gleichmaessig_kalibrierbar} ist dies genau dann möglich, wenn gemeinsame Schranken $L_{\max}$ und $M_{\max}$ existieren -- was für baugleiche Drohnen trivialerweise der Fall ist, da alle Drohnen dieselbe Dynamik besitzen. Das Schema $(\delta^*, \Delta t^*)$ kann damit einmal für den Drohnentyp berechnet und auf alle $k$ Drohnen des Schwarms angewendet werden, ohne die Kalibrierung $k$-fach durchzuführen.

\subsection{Zusammenfassung der Ergebnisse}

Dieses Kapitel hat eine systematische Theorie der Diskretisierung hybrider Systeme für die formale Verifikation entwickelt. Die zentralen Ergebnisse lassen sich wie folgt zusammenfassen:

\textbf{Fehlercharakterisierung.} Der Gesamtabstraktionsfehler setzt sich aus einem räumlichen Anteil $\frac{\sqrt{n}\delta}{2}$ und einem dynamischen Anteil $\frac{M\Delta t}{2L}(e^{LT}-1)$ zusammen (Theorem~\ref{thm:gesamtfehler}). Beide Anteile sind formal begrenzt und können durch Wahl der Parameter $\delta$ und $\Delta t$ kontrolliert werden.

\textbf{Universelle Kalibrierbarkeit.} Für jede Klasse von Systemen mit gemeinsamen Schranken $L_{\max}$ und $M_{\max}$ existiert ein einziges $\varepsilon$-kalibriertes Schema, das für alle Systeme der Klasse gleichzeitig gilt (Theorem~\ref{thm:gleichmaessig_kalibrierbar}). Dies ermöglicht klassenweite Zertifizierung ohne individuelle Systemanalyse.

\textbf{Sound Abstraktion.} Die $\varepsilon$-Erweiterung der Transitionsrelation gewährleistet Soundness: Safety-Beweise im diskreten Modell übertragen sich auf das reale hybride System (Theorem~\ref{thm:soundness}).

\textbf{Grenzen der Fehlervermeidung.} Vollständige Fehlervermeidung bei endlichem Zustandsraum ist genau für Systeme mit linearer Dynamik und kompatiblen Guards möglich (Theorem~\ref{thm:notwendig}). Für nichtlineare Systeme verbleibt stets ein positiver Restfehler, der jedoch durch Kalibrierung beliebig klein gemacht werden kann.

\textbf{Praktische Anwendbarkeit.} Für alle untersuchten konkreten Systeme -- TurtleBot3, autonomes Fahrzeug, KUKA-Roboter, Drohnen-Schwarm -- wurde das $\varepsilon$-kalibrierte Schema explizit berechnet und seine Praxistauglichkeit demonstriert. Die Auflösungen liegen im Bereich von 3--50\,mm, was mit moderner Sensorik und Steuerungshardware erreichbar ist.


\section{Erweiterte Themen}

Die in den vorangegangenen Kapiteln entwickelten Grundlagen -- temporale Logiken, PDDL Planning, der Subgraph Algorithmus und hybride Verifikation -- bilden das Fundament für eine Reihe weiterführender Themen, die in der Forschung aktiv bearbeitet werden und die Ausdrucksstärke sowie Anwendbarkeit der vorgestellten Verfahren erheblich erweitern. Dieses Kapitel behandelt drei solcher Erweiterungen ausführlich: temporale PDDL (PDDL 2.1) für zeitkritische Planungsprobleme, probabilistisches Planning für unsichere Umgebungen sowie hierarchische Planung (HTN) für strukturierte Aufgabendekompositionen. Für jede Erweiterung wird die formale Semantik vollständig entwickelt, die Verbindung zum Subgraph Algorithmus und temporalen Model Checking hergestellt und an konkreten Beispielen illustriert.

\subsection{Temporale PDDL (PDDL 2.1)}

Klassisches PDDL behandelt Aktionen als instantan: Eine Aktion wird angewendet, der Zustand ändert sich, und die nächste Aktion beginnt. In der Praxis haben Aktionen Dauern -- ein Roboter benötigt Zeit zum Fahren, ein Kran Zeit zum Heben, ein Prozess Zeit zum Abkühlen. PDDL 2.1 \cite{fox2003} erweitert PDDL um \emph{durative actions}, die eine explizite zeitliche Ausdehnung besitzen und deren Vor- und Nachbedingungen an verschiedenen Zeitpunkten der Aktionsausführung geprüft bzw.\ gesetzt werden.

\subsubsection{Formale Semantik durativer Aktionen}

\begin{definition}[Durative Aktion]
Eine durative Aktion ist ein 4-Tupel
\[ a = (\mathit{par},\; \mathit{dur},\; \mathit{cond},\; \mathit{eff}) \]
mit folgenden Komponenten:
\begin{itemize}
  \item[-] $\mathit{par}$: Parameterliste (wie in klassischem PDDL)
  \item[-] $\mathit{dur}$: Dauerconstraint, typischerweise $= d$ für eine feste Dauer $d \in \mathbb{R}_{>0}$
  \item[-] $\mathit{cond}$: zeitlich qualifizierte Bedingung, zusammengesetzt aus
    \begin{itemize}
      \item $\mathbf{at\;start}(\varphi)$: $\varphi$ gilt zum Startzeitpunkt $t_s$
      \item $\mathbf{over\;all}(\varphi)$: $\varphi$ gilt für alle $t \in [t_s, t_s + d]$
      \item $\mathbf{at\;end}(\varphi)$: $\varphi$ gilt zum Endzeitpunkt $t_s + d$
    \end{itemize}
  \item[-] $\mathit{eff}$: zeitlich qualifizierter Effekt, analog zu $\mathit{cond}$
\end{itemize}
\end{definition}

\begin{definition}[Temporaler Plan]
Ein temporaler Plan ist eine Menge von zeitgestempelten Aktionsinstanzen $\pi = \{(a_1, t_1), (a_2, t_2), \ldots, (a_n, t_n)\}$ mit $t_i \in \mathbb{R}_{\geq 0}$. Anders als im klassischen Plan können Aktionen überlappen: Aktionen $a_i$ und $a_j$ überlappen, wenn $[t_i, t_i + d_i] \cap [t_j, t_j + d_j] \neq \emptyset$.
\end{definition}

Die Semantik eines temporalen Plans wird durch einen zeitlichen Zustandsbeobachter definiert, der zu jedem Zeitpunkt $t$ den aktuellen Zustand $s(t)$ bestimmt. Dabei entstehen Konflikte, wenn zwei gleichzeitig aktive Aktionen auf dasselbe Prädikat zugreifen:

\begin{definition}[Mutex-Konflikt zwischen durativen Aktionen]
Zwei durative Aktionen $a_i$ und $a_j$ sind in einem temporalen Plan \emph{mutex}, wenn ihre Ausführungszeitfenster sich überlappen und mindestens eine der folgenden Bedingungen gilt:
\begin{enumerate}
  \item $a_i$ löscht ein Prädikat, das $a_j$ als \texttt{over-all}-Bedingung benötigt.
  \item $a_i$ und $a_j$ setzen dasselbe Prädikat auf verschiedene Werte.
  \item $a_i$ benötigt ein Prädikat, das $a_j$ zu einem Zeitpunkt innerhalb der Überlappung löscht.
\end{enumerate}
\end{definition}

\subsubsection{PDDL 2.1 Beispiel: Logistikdomäne mit Fahrtzeiten}

\begin{lstlisting}[language=PDDL, caption=Durative Action in PDDL 2.1 -- Logistik mit Fahrtdauer]
(:durative-action drive-truck
  :parameters (?truck - truck ?from ?to - location ?city - city)
  :duration (= ?duration (travel-time ?from ?to))
  :condition (and
    (at start (at ?truck ?from))
    (at start (in-city ?from ?city))
    (at start (in-city ?to ?city))
    (over all (road-open ?from ?to)))
  :effect (and
    (at start  (not (at ?truck ?from)))
    (at end    (at ?truck ?to))))
\end{lstlisting}

Das Schlüsselwort \texttt{(travel-time ?from ?to)} bezeichnet eine numerische Funktion, die die Fahrtdauer zwischen zwei Orten aus einer Nachschlagetabelle liefert. Die Bedingung \texttt{(over all (road-open ?from ?to))} stellt sicher, dass die Straße während der gesamten Fahrt befahrbar bleibt -- ein Constraint, das in klassischem PDDL nicht ausdrückbar ist.

\subsubsection{Temporale Verifikation mit CTL}

Die temporale Verifikation von PDDL-2.1-Plänen erfordert eine Erweiterung der Kripke-Struktur um eine Zeitdimension. Wir definieren eine \emph{zeiterweiterte Kripke-Struktur}:

\begin{definition}[Zeiterweiterte Kripke-Struktur]
Eine zeiterweiterte Kripke-Struktur ist ein Tupel $M_T = (S, T, R, L)$ mit:
\begin{itemize}
  \item[-] $S$: Zustandsmenge (wie zuvor)
  \item[-] $T \subseteq \mathbb{R}_{\geq 0}$: diskretisierte Zeitpunktmenge
  \item[-] $R \subseteq (S \times T) \times (S \times T)$: zeitbehaftete Transitionsrelation
  \item[-] $L: S \times T \to 2^{AP}$: zeitabhängige Labeling-Funktion
\end{itemize}
\end{definition}

Auf dieser Struktur lassen sich zeitkritische CTL-Eigenschaften spezifizieren, etwa:
\begin{align}
  \mathbf{AG}_{t \leq T_{\max}}(\mathit{goal\text{-}reachable}) \quad &\text{(Ziel innerhalb Zeitfenster erreichbar)} \\
  \mathbf{AG}(\mathit{deadline}(a_i) \to \mathbf{AF}_{\leq d_i} \mathit{complete}(a_i)) \\
  &\quad\text{(jede Aktion schließt in ihrer Dauer ab)}
\end{align}

\subsubsection{Erweiterung des Subgraph Algorithmus für PDDL 2.1}

Für zeiterweiterte Zustände wird die Signatur um eine Zeitkomponente ergänzt. Die zeiterweiterte Signatur $\sigma_j^T$ eines Zustands $(s, t)$ lautet:
\begin{align}
  \sigma_j^T(s, t) = \sigma_j(s) + \lfloor t / \Delta t \rfloor \cdot 2^{2n}
\end{align}
wobei $\Delta t$ die Zeitauflösung und $\sigma_j(s)$ die klassische Signatur ist. Diese Erweiterung bleibt eindeutig (da die Zeitkomponente in einem disjunkten Zahlenbereich liegt) und ermöglicht zeitkritisches Pattern-Matching in $O(n^3 \cdot |T|)$, also linear in der Anzahl der Zeitschritte.

\begin{theorem}[Korrektheit zeiterweiterter Signaturen]
Die zeiterweiterte Signatur $\sigma^T$ ist injektiv auf $(S \times T)$: Verschiedene zeitgestempelte Zustände erhalten verschiedene Signaturen.
\end{theorem}

\begin{proof}
Seien $(s_1, t_1) \neq (s_2, t_2)$. Fall 1: $s_1 \neq s_2$. Dann unterscheiden sich bereits die klassischen Anteile $\sigma_j(s_1) \neq \sigma_j(s_2)$ (nach Lemma 7.1), also $\sigma_j^T(s_1, t_1) \neq \sigma_j^T(s_2, t_2)$. Fall 2: $s_1 = s_2$, $t_1 \neq t_2$. Dann $\lfloor t_1/\Delta t\rfloor \neq \lfloor t_2/\Delta t\rfloor$ (für $|t_1 - t_2| \geq \Delta t$), und da der Zeitterm $\lfloor t/\Delta t\rfloor \cdot 2^{2n}$ in einem Bereich liegt, der disjunkt vom klassischen Signaturbereich $[0, n \cdot 2^n)$ ist, folgt die Eindeutigkeit.
\end{proof}

\subsubsection{Komplexitätsanalyse temporaler Pläne}

\begin{theorem}[Komplexität temporaler Planvalidierung]
Die Validierung eines temporalen Plans $\pi$ mit $n$ Aktionen und Zeitfenster $[0, T_{\max}]$ gegen eine Menge von $k$ CTL-Eigenschaften hat mit dem zeiterweiterten Subgraph Algorithmus Komplexität:
\begin{align}
  O\!\left(n \cdot \left\lceil \frac{T_{\max}}{\Delta t} \right\rceil \cdot m^3 \cdot k\right)
\end{align}
wobei $m$ die Größe der größten Property-Pattern ist.
\end{theorem}

\begin{proof}
Für jeden der $n$ Planschritte und jeden der $\lceil T_{\max}/\Delta t \rceil$ Zeitpunkte wird ein zeitgestempelter Zustand erzeugt. Für jede der $k$ Eigenschaften wird der zeiterweiterte Subgraph Algorithmus mit Komplexität $O(m^3)$ angewendet. Die Gesamtkomplexität ergibt sich durch Multiplikation.
\end{proof}

\subsection{Probabilistisches Planning}

Klassisches PDDL setzt deterministische Aktionen voraus: Jede Aktion hat genau einen Effekt. In der Realität sind Aktionen jedoch häufig unsicher -- ein Greifarm kann einen Gegenstand verfehlen, ein Sensor kann fehlerhafte Messwerte liefern, ein Kommunikationskanal kann ausfallen. Probabilistisches Planning modelliert diese Unsicherheit explizit durch Wahrscheinlichkeitsverteilungen über Aktionseffekte.

\subsubsection{Formale Grundlagen: Markov-Entscheidungsprozesse}

\begin{definition}[Markov-Entscheidungsprozess (MDP)]
Ein MDP ist ein Tupel $\mathcal{M} = (S, A, P, R, \gamma)$ mit:
\begin{itemize}
  \item[-] $S$: endliche Zustandsmenge
  \item[-] $A$: endliche Aktionsmenge
  \item[-] $P: S \times A \times S \to [0,1]$: Ãœbergangswahrscheinlichkeit, wobei $\sum_{s'} P(s,a,s') = 1$
  \item[-] $R: S \times A \to \mathbb{R}$: Belohnungsfunktion
  \item[-] $\gamma \in [0,1)$: Diskontierungsfaktor
\end{itemize}
\end{definition}

Jeder probabilistische PDDL-Plan induziert
einen MDP; Komponenten entsprechen den PDDL-Konstrukten.

\subsubsection{Probabilistisches PDDL}

Das folgende Beispiel zeigt einen unzuverlässigen Greifarm, der einen Ball mit Wahrscheinlichkeit 0,9 erfolgreich aufnimmt:

\begin{lstlisting}[language=PDDL, caption=Probabilistischer PDDL-Operator]
(:action pick-unreliable
  :parameters (?b - ball ?r - room ?g - gripper)
  :precondition (and (at-ball ?b ?r) (at-robby ?r) (free ?g))
  :effect (probabilistic
    0.90 (and (carry ?b ?g)
              (not (free ?g))
              (not (at-ball ?b ?r)))
    0.07 (and (damaged ?b)
              (free ?g))
    0.03 (free ?g)))
\end{lstlisting}

Mit Wahrscheinlichkeit 0,90 wird der Ball erfolgreich gegriffen, mit 0,07 wird er beim Greifversuch beschädigt (und bleibt im Raum), und mit 0,03 schlägt der Greifversuch ohne Beschädigung fehl.

\subsubsection{Probabilistic CTL (PCTL)}

Für die Verifikation probabilistischer Systeme wird CTL zur \emph{Probabilistic CTL} (PCTL) erweitert \cite{baier2008}:

\begin{definition}[PCTL-Syntax]
PCTL-Formeln sind induktiv definiert durch:
\begin{align}
  \Phi ::= \text{true} \mid p \mid \neg\Phi \mid \Phi_1 \land \Phi_2 \mid \mathcal{P}_{\bowtie q}[\psi]
\end{align}
wobei $\psi$ ein Pfadformeln-Operator ist:
\begin{align}
  \psi ::= \mathbf{X}\Phi \mid \Phi_1 \mathbf{U}^{\leq k} \Phi_2 \mid \Phi_1 \mathbf{U} \Phi_2
\end{align}
Dabei ist $\mathcal{P}_{\bowtie q}[\psi]$ mit $\bowtie\, \in \{<, \leq, >, \geq\}$ und $q \in [0,1]$ eine Wahrscheinlichkeitsaussage: ``Die Wahrscheinlichkeit, dass $\psi$ gilt, ist $\bowtie q$.''
\end{definition}

\begin{beispiel}[PCTL-Formeln für Gripper-Domäne]
Für $p = 0{,}9$ sind folgende PCTL-Formeln relevant:
\begin{enumerate}
  \item \textit{Mit Wahrscheinlichkeit $\geq 0{,}99$ werden alle Bälle schließlich transportiert:}
  \begin{align}
    \mathcal{P}_{\geq 0{,}99}[\mathbf{F}(\mathit{all\text{-}delivered})]
  \end{align}
  \item \textit{Mit Wahrscheinlichkeit $< 0{,}05$ wird ein Ball beschädigt:}
  \begin{align}
    \mathcal{P}_{< 0{,}05}[\mathbf{F}(\mathit{damaged}(b))]
  \end{align}
  \item \textit{Innerhalb von 20 Schritten wird das Ziel mit Wahrscheinlichkeit $\geq 0{,}95$ erreicht:}
  \begin{align}
    \mathcal{P}_{\geq 0{,}95}[\mathbf{F}^{\leq 20}(\mathit{goal})]
  \end{align}
\end{enumerate}
\end{beispiel}

\subsubsection{Wertiteration für MDPs}

Die optimale Strategie (Policy) für einen MDP minimiert oder maximiert den erwarteten kumulierten Nutzen. Der klassische Algorithmus ist die \emph{Wertiteration}:

\begin{algorithm}[H]
  \caption{Wertiteration für MDPs}
  \begin{algorithmic}[1]
    \Eingabe MDP $\mathcal{M} = (S, A, P, R, \gamma)$, Konvergenzkriterium $\varepsilon$
    \Ausgabe Optimale Wertefunktion $V^*$, optimale Policy $\pi^*$
    \Prozedur{Wertiteration}{$\mathcal{M}, \varepsilon$}
      \State $V_0(s) \gets 0$ für alle $s \in S$
      \State $k \gets 0$
      \Wiederhole
        \FuerAlle{$s \in S$}
          \State $V_{k+1}(s) \gets \max_{a \in A} \left[ R(s,a) + \gamma \sum_{s'} P(s,a,s') V_k(s') \right]$
        \EndeFuer
        \State $k \gets k + 1$
      \Bis{$\max_{s \in S} |V_k(s) - V_{k-1}(s)| < \varepsilon$}
      \State $\pi^*(s) \gets \arg\max_{a \in A} \left[ R(s,a) + \gamma \sum_{s'} P(s,a,s') V^*(s') \right]$
      \State \Return $(V^*, \pi^*)$
    \EndProzedur
  \end{algorithmic}
\end{algorithm}

\begin{theorem}[Konvergenz der Wertiteration]
Für einen endlichen MDP mit Diskontierungsfaktor $\gamma \in [0,1)$ konvergiert die Wertiteration gegen die eindeutige optimale Wertefunktion $V^*$. Die Konvergenzrate ist geometrisch mit Faktor $\gamma$: Nach $k$ Iterationen gilt
\begin{align}
  \|V_k - V^*\|_\infty \leq \frac{\gamma^k}{1-\gamma} \|V_1 - V_0\|_\infty.
\end{align}
\end{theorem}

\begin{proof}
Der Bellman-Operator $\mathcal{T}$ ist eine $\gamma$-Kontraktion bzgl.\ $\|\cdot\|_\infty$. Es gilt:
\begin{align}
  \|(\mathcal{T} V) - (\mathcal{T} W)\|_\infty &= \max_s \left|\max_a \gamma \sum_{s'} P(s,a,s')(V(s') - W(s'))\right| \\
  &\leq \gamma \max_s \max_a \sum_{s'} P(s,a,s') |V(s') - W(s')| \notag \\
  &\leq \gamma \|V - W\|_\infty. \notag
\end{align}
Nach dem Banachschen Fixpunktsatz besitzt jede Kontraktion auf einem vollständigen metrischen Raum genau einen Fixpunkt, und die Iterationsfolge $V_k = \mathcal{T}^k V_0$ konvergiert geometrisch gegen diesen. Die Fehlerabschätzung ergibt sich durch geometrische Summation.
\end{proof}

\subsubsection{Probabilistischer Subgraph Algorithmus}

Für probabilistische Systeme wird die Signatur um eine Wahrscheinlichkeitskomponente erweitert. Statt einer Adjazenzmatrix $A \in \{0,1\}^{n \times n}$ verwendet man eine Wahrscheinlichkeitsmatrix $P \in [0,1]^{n \times n}$, wobei $P[i,j]$ die Übergangswahrscheinlichkeit von Zustand $i$ nach $j$ angibt. Die probabilistische Signatur lautet:
\begin{align}
  \sigma_j^P(s) = \sum_{i=0}^{n-1} \lfloor P[i,j] \cdot 2^b \rfloor \cdot 2^{ib} + j \cdot 2^{nb}
\end{align}
wobei $b$ die Anzahl der Bits für die Quantisierung der Wahrscheinlichkeit ist (typischerweise $b = 8$ für 256 Stufen). Das Matching verwendet dann eine tolerante Variante der LCS, die Wahrscheinlichkeitsdifferenzen bis zu einem Schwellwert $\tau$ akzeptiert.

\begin{definition}[Tolerantes Wahrscheinlichkeits-Matching]
Zwei probabilistische Signaturen $\sigma^P(s_1)$ und $\sigma^P(s_2)$ \emph{$\tau$-matchen}, wenn für alle Komponenten $j$ gilt:
\begin{align}
  |\sigma_j^P(s_1) - \sigma_j^P(s_2)| \leq \tau \cdot 2^{nb}.
\end{align}
\end{definition}

\subsubsection{Anwendung: Verifikation des Gripper-Plans unter Unsicherheit}

Für den probabilistischen Gripper-Plan mit Greifwahrscheinlichkeit $p = 0{,}9$ und drei Bällen lässt sich die Erfolgswahrscheinlichkeit für den vollständigen Transport berechnen. Da jede Übergabe zwei Greifvorgänge erfordert (Ball greifen + Ball ablegen), und der Transport mit dem optimalen Plan in 3 Fahrten mit je 2 Greifvorgängen durchgeführt wird:
\begin{align}
  P(\text{Erfolg}) = p^6 = 0{,}9^6 \approx 0{,}531.
\end{align}
Dies ist unbefriedigend niedrig. Mit Wiederholungsversuchen bei Fehlschlag verbessert sich die Erfolgswahrscheinlichkeit pro Greifvorgang auf:
\begin{align}
  P(\text{Erfolg nach max.\ $k$ Versuchen}) = 1 - (1-p)^k.
\end{align}
Für $k = 3$ Versuche: $1 - 0{,}1^3 = 0{,}999$, und der Gesamtplan erreicht $0{,}999^6 \approx 0{,}994$, was die PCTL-Spezifikation $\mathcal{P}_{\geq 0{,}99}[\mathbf{F}(\mathit{all\text{-}delivered})]$ erfüllt.

\subsection{Hierarchische Planung (HTN)}

Hierarchisches Task Network Planning (HTN) \cite{erol1994} ist ein Planungsparadigma, das komplexe Aufgaben durch schrittweise Dekomposition in einfachere Teilaufgaben löst. HTN-Planer imitieren das menschliche Vorgehen: Man beginnt mit einer abstrakten Aufgabe (``liefere Paket''), zerlegt sie in Teilschritte (``lade Paket, fahre, entlade Paket'') und verfeinert diese weiter, bis primitive, direkt ausführbare Aktionen entstehen.

\subsubsection{Formale Definition: HTN-Planungsproblem}

\begin{definition}[HTN-Planungsproblem]
Ein HTN-Planungsproblem ist ein Tupel $\Pi = (s_0, T_0, D, \mathcal{M})$ mit:
\begin{itemize}
  \item[-] $s_0 \in S$: Initialzustand
  \item[-] $T_0$: initiale Taskliste (geordnet oder ungeordnet)
  \item[-] $D$: Menge primitiver Aktionen (wie in klassischem PDDL)
  \item[-] $\mathcal{M}$: Menge von Dekompositionsmethoden
\end{itemize}
Eine Dekompositionsmethode $m \in \mathcal{M}$ ist ein Tupel $m = (\mathit{task}, \mathit{pre}, \mathit{subtasks})$ mit:
\begin{itemize}
  \item[-] $\mathit{task}$: die abstrakte Aufgabe, die $m$ dekomponiert
  \item[-] $\mathit{pre}$: Vorbedingung, die für die Anwendbarkeit von $m$ gelten muss
  \item[-] $\mathit{subtasks}$: geordnete oder ungeordnete Liste von Unteraufgaben (primitiv oder abstrakt)
\end{itemize}
\end{definition}

Eine Lösung für $\Pi$ ist ein primitiver Plan $\pi$ (eine Sequenz primitiver Aktionen), der durch sukzessive Anwendung von Dekompositionsmethoden aus $T_0$ entsteht und den Initialzustand in einen Zielzustand überführt.

\subsubsection{HTN-Beispiel: Transport-Domäne}

Die folgende HTN-Domäne beschreibt den Transport von Paketen mit mehreren Fahrzeugen. Die abstrakte Aufgabe \texttt{transport-package} wird je nach Distanz in zwei verschiedene Methoden dekomponiert:

\begin{lstlisting}[language=PDDL, caption=HTN-Methoden für Transport-Domäne]
; Methode 1: Lokaler Transport (innerhalb derselben Stadt)
(:method local-transport
  :parameters (?pkg - package ?from ?to - location)
  :task (transport-package ?pkg ?from ?to)
  :precondition (same-city ?from ?to)
  :subtasks (and
    (task0 (find-truck ?from))
    (task1 (load-package ?pkg ?from))
    (task2 (drive-to-destination ?to))
    (task3 (unload-package ?pkg ?to))
    :ordering (< task0 task1)
             (< task1 task2)
             (< task2 task3)))

; Methode 2: Ueberregionaler Transport (via Flughafen)
(:method air-transport
  :parameters (?pkg - package ?from ?to - location
               ?hub-from ?hub-to - airport)
  :task (transport-package ?pkg ?from ?to)
  :precondition (not (same-city ?from ?to))
  :subtasks (and
    (task0 (transport-package ?pkg ?from ?hub-from))
    (task1 (load-airplane ?pkg ?hub-from))
    (task2 (fly ?hub-from ?hub-to))
    (task3 (unload-airplane ?pkg ?hub-to))
    (task4 (transport-package ?pkg ?hub-to ?to))
    :ordering (< task0 task1) (< task1 task2)
             (< task2 task3) (< task3 task4)))
\end{lstlisting}

\subsubsection{HTN-Suchraum und Komplexität}

\begin{definition}[Decomposition Tree]
Ein Dekompositionsbaum für ein HTN-Problem ist ein geordneter Baum, dessen Wurzel die initiale Taskliste $T_0$ darstellt, dessen innere Knoten abstrakte Aufgaben sind, die durch Dekompositionsmethoden zerlegt werden, und dessen Blätter primitive Aktionen sind. Der Plan ergibt sich als die geordnete Sequenz der Blätter.
\end{definition}

\begin{theorem}[Komplexität des HTN-Planens]
HTN-Planen ist EXPTIME-vollständig. Vollständig geordnete Tasklisten (SO-HTN) sind PSPACE-vollständig; einfache HTN-Probleme ohne Zustandsbedingungen NP-vollständig.
\end{theorem}

Die Komplexitätshierarchie zeigt, dass Strukturierungswissen die Planungskomplexität erheblich beeinflussen kann: Ein vollständig geordnetes HTN ist effektiver suchbar als ein ungeordnetes, weil die Methoden stärkere Vorauswahl treffen.

\subsubsection{Verbindung zu temporalem Model Checking}

HTN-Pläne lassen sich als Kripke-Strukturen modellieren, wobei die Hierarchieebenen explizit dargestellt werden:

\begin{definition}[HTN-Kripke-Struktur]
Für ein HTN-Planungsproblem $\Pi$ definieren wir die HTN-Kripke-Struktur $M_{\mathit{HTN}} = (S_H, R_H, L_H)$ mit:
\begin{itemize}
  \item[-] $S_H$: Menge partiell
  dekomponierten Tasklisten (alle Zwischenzustände)
  \item[-] $R_H$: Ãœbergang durch Anwendung einer Dekompositionsmethode oder primitiven Aktion
  \item[-] $L_H$: Labels für Abstrktionsebene, aktuelle Tasks, und Planzustand
\end{itemize}
\end{definition}

Auf $M_{\mathit{HTN}}$ lassen sich nun strukturbezogene Eigenschaften formal spezifizieren:

\begin{beispiel}[HTN-Verifikationseigenschaften]
Für die Transport-HTN-Domäne seien folgende Eigenschaften zu verifizieren:
\begin{enumerate}
  \item \textit{Korrektheit der Dekomposition}: Jede abstrakte Aufgabe wird schließlich vollständig in primitive Aktionen zerlegt:
  \begin{align}
    \mathbf{AG}(\mathit{abstract\text{-}task}(T) \to \mathbf{AF}(\mathit{decomposed}(T)))
  \end{align}
  \item \textit{Abhängigkeitsfreiheit}: Keine zirkuläre Abhängigkeit zwischen Teilaufgaben:
  \begin{align}
    \mathbf{AG}\neg(\mathit{subtask\text{-}of}(T_1, T_2) \land \mathit{subtask\text{-}of}(T_2, T_1))
  \end{align}
  \item \textit{Zeitkonsistenz}: In der SO-HTN werden Teilaufgaben in der vorgeschriebenen Reihenfolge ausgeführt:
  \begin{align}
    \mathbf{AG}(\mathit{order}(T_1 < T_2) \to \mathbf{AF}(\mathit{complete}(T_1) \land \mathbf{AX}\mathit{start}(T_2)))
  \end{align}
\end{enumerate}
\end{beispiel}

\subsubsection{Subgraph-basierte Methodenauswahl}

Ein praktisch wichtiges Problem im HTN-Planen ist die effiziente Auswahl der anwendbaren Dekompositionsmethoden in einem gegebenen Zustand. Der Subgraph Algorithmus kann hier eingesetzt werden, um den aktuellen Zustand mit den Vorbedingungsmustern der Methoden zu vergleichen:

\begin{algorithm}[H]
  \caption{Subgraph-basierte HTN-Methodenauswahl}
  \begin{algorithmic}[1]
    \Eingabe Aktueller Zustand $s$, Taskliste $T$, Methodenmenge $\mathcal{M}$
    \Ausgabe Menge anwendbarer Methoden $\mathcal{M}_{\mathit{app}}$
    \Prozedur{MethodenAuswahl}{$s, T, \mathcal{M}$}
      \State $\sigma_s \gets \textsc{BerechneSig}(s)$ \Comment{Einmalig, $O(n^2)$}
      \State $\mathcal{M}_{\mathit{app}} \gets \emptyset$
      \FuerAlle{$m \in \mathcal{M}$ mit $m.\mathit{task} \in T$}
        \State $\sigma_m \gets \textsc{BerechneSig}(m.\mathit{pre})$
        \Wenn{\textsc{SubgraphMatch}($\sigma_m$, $\sigma_s$)}
          \State $\mathcal{M}_{\mathit{app}} \gets \mathcal{M}_{\mathit{app}} \cup \{m\}$
        \EndeWenn
      \EndeFuer
      \State \Return $\mathcal{M}_{\mathit{app}}$
    \EndProzedur
  \end{algorithmic}
\end{algorithm}

{\sloppy\begin{theorem}[Korrektheit der Subgraph-basierten Methodenauswahl]
Algorithmus~12.2 ist \emph{sound}; bei vollständigem Matching auch \emph{vollständig}.
\end{theorem}}\par

\begin{proof}
\textit{Soundness}: Wenn \textsc{SubgraphMatch}($\sigma_m$, $\sigma_s$) wahr zurückgibt, so existiert nach Definition des Subgraph-Matching eine Einbettung des Vorbedingungsmusters $m.\mathit{pre}$ in den aktuellen Zustand $s$, d.\,h., alle Prädikate der Vorbedingung sind in $s$ erfüllt. Die Methode ist damit korrekt anwendbar.

\textit{Vollständigkeit}: Falls der Subgraph Algorithmus unter Rotationsbeschränkung eine anwendbare Methode nicht erkennt (false negative), so liegt ein Vollständigkeitsproblem des Matching-Algorithmus vor, nicht der Methodenauswahl selbst. Für strukturierte Planungsprobleme, in denen die Vorbedingungsmuster zyklisch geordnet sind, ist die Rotationsbeschränkung vollständig (vgl.\ Abschnitt~7.3), und damit auch die Methodenauswahl.
\end{proof}

\subsubsection{Vergleich: Klassisches PDDL vs.\ PDDL 2.1 vs.\ HTN}

\begin{table}[h]
  \centering
  \caption{Vergleich der drei PDDL-Paradigmen}
  \label{tab:pddl_paradigmen}
  \begin{tabular}{p{2.4cm}p{2.8cm}p{3.0cm}p{2.8cm}}
    \toprule
    Merkmal & Klassisches PDDL & PDDL 2.1 (temporal) & HTN \\
    \midrule
    Aktionsdauer & Instantan & Explizit ($d \in \mathbb{R}$) & Methoden-abhängig \\
    Parallelität & Nein & Ja (mit Mutex) & Ja (ungeordnete Tasks) \\
    Struktur\-wissen & Nein & Nein & Ja (Methoden) \\
    Komplexität & PSPACE & PSPACE & EXPTIME \\
    Verifikation & CTL/LTL & TCTL & HTN-CTL \\
    Subgraph & $O(n^3)$ & $O(n^3 |T|)$ & $O(n^3 |\mathcal{M}|)$ \\
    \bottomrule
  \end{tabular}
\end{table}

Die Tabelle zeigt, dass alle drei Paradigmen mit dem Subgraph Algorithmus kompatibel sind: Die zeitliche bzw.\ hierarchische Erweiterung erhöht die Komplexität lediglich linear in der Anzahl der Zeitschritte bzw.\ Methoden, was den polynomialen Charakter des Gesamtansatzes bewahrt.

\subsection{Zusammenfassung}

Dieses Kapitel hat drei wesentliche Erweiterungen des PDDL-Formalismus ausführlich behandelt. Temporales PDDL (PDDL 2.1) ermöglicht die explizite Modellierung von Aktionsdauern, paralleler Aktionsausführung und zeitkritischen Constraints. Die formale Semantik durativer Aktionen wurde vollständig entwickelt, und die Erweiterung des Subgraph Algorithmus auf zeiterweiterte Signaturen wurde mit einem Korrektheitsbeweis versehen. Probabilistisches Planning modelliert Unsicherheit durch MDPs und PCTL-Spezifikationen. Der Konvergenzbeweis der Wertiteration stellt die theoretische Basis für die Berechnung optimaler Policies unter Unsicherheit bereit, und die probabilistische Signaturerweiterung integriert diese in den Subgraph-Rahmen. HTN-Planung schließlich nutzt hierarchisches Strukturwissen, um komplexe Planungsaufgaben durch Dekomposition handhabbar zu machen. Die formale HTN-Kripke-Struktur ermöglicht die Verifikation strukturbezogener Eigenschaften, und die Subgraph-basierte Methodenauswahl beschleunigt die HTN-Suche effizient. In allen drei Fällen bleibt der Subgraph Algorithmus der zentrale Verifikationsbaustein, der die Komplexität strukturierter Planungsprobleme beherrschbar macht.

\section{Schluss}

\subsection{Hauptergebnisse}

Diese Arbeit hat temporales Model Checking und PDDL-Planungssysteme behandelt. Als zentrale Innovation wurde der Subgraph Algorithmus eingeführt.

\textbf{Temporale Logiken:} LTL und CTL bieten mächtige Formalismen zur Spezifikation zeitabhängiger Eigenschaften. CTL-Modellprüfung ist in polynomieller Zeit lösbar; LTL-Modellprüfung ist PSPACE-vollständig. Bei großen Zustandsräumen stellen beide Komplexitäten erhebliche Herausforderungen dar -- der Subgraph Algorithmus begegnet genau diesem Problem.

\textbf{PDDL als Standardformalismus:} Vier klassische Domänen -- Blocksworld, Logistics, Gripper und weitere Varianten -- wurden mit vollständigen Beispielen präsentiert. Planungsalgorithmen von der Forward-~und Backward-Suche über GraphPlan und FF bis hin zu SAT-basierter Planung wurden vorgestellt und verglichen.

\textbf{Der Subgraph Algorithmus als zentrale Innovation:} Kapitel~7 stellt den Hauptbeitrag dieser Arbeit dar: einen signaturbasierten Algorithmus, der das State Explosion Problem durch strukturelle Abstraktion überwindet.

\begin{table}[h]
	\centering
	\caption{Vergleich mit etablierten Ansätzen}
	\label{tab:final_comparison}
	\begin{tabular}{llll}
		\toprule
		Ansatz & Komplexität & Vollständigkeit & Speedup \\
		\midrule
		SPIN (explizit) & $O(2^n)$ & Ja & Baseline \\
		NuSMV (symbolisch) & Variabel (BDD) & Ja & $2$--$10\times$ \\
		VF2 (Subgraph Iso.) & $O(n!)$ & Ja & -- \\
		\textbf{Subgraph-MC} & \textbf{$O(n^3)$} & \textbf{Pattern-basiert} & \textbf{$15$--$1239\times$} \\
		\bottomrule
	\end{tabular}
\end{table}

Die wesentlichen Eigenschaften des Subgraph Algorithmus:
\begin{itemize}
	\item[-] \textbf{Polynomiale Komplexität}: Reduktion von $O(n!)$ auf $O(n^3)$ für Pattern-Matching
	\item[-] \textbf{Signaturbasierte Kodierung}: Kompakte, eindeutige Repräsentation von Zuständen
	\item[-] \textbf{Rotationsbasierte Suche}: Effiziente Behandlung von Objektpermutationen
	\item[-] \textbf{Praktische Effizienz}: $15$--$1239\times$ Speedup in experimenteller Evaluierung
\end{itemize}

\textbf{Integration von Planung und Verifikation:} Kapitel~9 und~10 haben gezeigt, wie klassische Planungsalgorithmen mit dem Subgraph-basierten Model Checking kombiniert werden. Der vollständige Workflow umfasst: Plangenerierung, Signaturextraktion, Pattern-Verifikation und formale Zertifizierung -- alles mit polynomialem Gesamtaufwand für den Verifikationsschritt.

\subsection{Wissenschaftlicher Beitrag}

Die Arbeit leistet folgende wissenschaftliche Beiträge:

\begin{enumerate}
	\item \textbf{Theoretische Fundierung}: Formale Definition der signaturbasierten Zustandsrepräsentation mit Beweis der Eindeutigkeit (Lemma~7.1) und Korrektheit
	\item \textbf{Algorithmische Innovation}: $O(n^3)$ Algorithmus für Pattern-Matching, wo klassische Ansätze exponentielle Komplexität haben
	\item \textbf{Integration in Model Checking}: Einbettung in LTL-~und CTL-Workflows
	      sowie in klassische Planungsalgorithmen
	\item \textbf{Experimentelle Validierung}: Evaluierung mit drei Domänen,
	      die $15$--$1239\times$ Beschleunigung demonstriert
	\item \textbf{Praktische Anwendbarkeit}: Nachweis für reale
	      PDDL-Probleme in Robotik, Logistik und autonomen Systemen
\end{enumerate}

\subsection{Limitationen}

Bei aller Leistungsfähigkeit hat der Ansatz bekannte Einschränkungen:

\textbf{Rotationsbeschränkung}: Die Einschränkung auf zyklische Rotationen kann in bestimmten Fällen mögliche Matches übersehen. Der Algorithmus ist \textit{sound} -- er liefert keine false positives --, ist jedoch nicht vollständig \textit{complete}. Dieser Trade-off wird bewusst zugunsten polynomialer Laufzeit in Kauf genommen.

\textbf{Strukturierte Probleme}: Bei Systemen ohne wiederkehrende Strukturmuster bietet der Subgraph Algorithmus weniger Vorteile. Er ist komplementär zu klassischen Ansätzen und optimal für strukturierte Planungsdomänen.

\textbf{Probabilistische Systeme}: Die aktuelle Formulierung setzt deterministisches Verhalten voraus. Erweiterungen auf PCTL erfordern angepasste Signaturen.

\subsection{Ausblick}

Für die Zukunft ergeben sich mehrere vielversprechende Forschungsrichtungen:

\textbf{Algorithmische Erweiterungen:}
\begin{itemize}
	\item[-] Kombination mit BDDs für sehr große Zustandsräume ($|S| > 10^5$)
	\item[-] GPU-Beschleunigung der parallelen Signatur-Berechnung
	\item[-] Erweiterung auf probabilistische Systeme (Signaturen für MDPs)
	\item[-] Inkrementelle Verifikation bei lokalen Planänderungen
\end{itemize}

\textbf{Tool-Integration:}
\begin{itemize}
	\item[-] Plugins für Fast Downward und PDDL4J
	\item[-] Anbindung an NuSMV und SPIN als Fallback
	      für komplexe temporale Formeln
	\item[-] Open-Source-Release mit ROS-Integration
	\item[-] Grafische Oberfläche zur Darstellung
	      von Signaturen und Matching-Ergebnissen
\end{itemize}

\textbf{Neue Anwendungsdomänen:}
\begin{itemize}
	\item[-] Cyber-Physical Systems und Smart Grids
	\item[-] Autonome Fahrzeuge mit Echtzeit-Pfadverifikation
	\item[-] Medizinische Entscheidungssysteme mit Safety-Garantien
	\item[-] Schwarmrobotik mit verteilter, kompositionaler Verifikation
\end{itemize}

\textbf{Theoretische Vertiefung:}
\begin{itemize}
	\item[-] Maschinell verifizierbare Beweise in Coq oder Isabelle/HOL
	\item[-] Formale Bisimulation zwischen kontinuierlichem und diskretem Modell (vertieft in Kapitel~11)
	\item[-] Komplexitätsuntergrenzen unter realistischen Strukturannahmen
	\item[-] Integration mit Machine Learning zur Verifikation gelernter Policies
\end{itemize}

Zusammenfassend lässt sich festhalten, dass temporales Model Checking, PDDL Planning und der Subgraph Algorithmus zusammen einen mächtigen, effizienten Ansatz für die Entwicklung und Verifikation intelligenter, sicherer und zuverlässiger Systeme bilden. Der Subgraph Algorithmus positioniert sich dabei als komplementäre Methode zu klassischen Verfahren: Für strukturelle Eigenschaften bietet er dramatische Effizienzgewinne, während klassisches CTL/LTL Model Checking für komplexe temporale Formeln weiterhin seine Stärken ausspielt. Die optimale Strategie kombiniert beide Ansätze situationsgerecht.

Die theoretischen Grundlagen sind solide, die Algorithmen sind effizient implementierbar, und die praktische Anwendbarkeit ist durch Experimente und Fallstudien belegt. Diese Arbeit legt das Fundament für eine neue Generation formal verifizierter Planungssysteme -- ein Schritt hin zu vertrauenswürdiger, beweisbar korrekter künstlicher Intelligenz.


\begin{thebibliography}{99}

\bibitem{baier2008}
Baier, C., Katoen, J.-P.:
\textit{Principles of Model Checking}.
MIT Press, Cambridge, MA, 2008.

\bibitem{clarke1999}
Clarke, E.\,M., Grumberg, O., Peled, D.\,A.:
\textit{Model Checking}.
MIT Press, Cambridge, MA, 1999.

\bibitem{clarke2018}
Clarke, E.\,M., Henzinger, T.\,A., Veith, H., Bloem, R. (Hrsg.):
\textit{Handbook of Model Checking}.
Springer, Cham, 2018.

\bibitem{pnueli1977}
Pnueli, A.:
The temporal logic of programs.
In: \textit{Proceedings of the 18th Annual Symposium on Foundations of Computer Science (FOCS)}, S.~46--57, IEEE, 1977.

\bibitem{emerson1990}
Emerson, E.\,A.:
Temporal and modal logic.
In: van Leeuwen, J. (Hrsg.), \textit{Handbook of Theoretical Computer Science}, Bd.~B, S.~995--1072. Elsevier, Amsterdam, 1990.

\bibitem{queille1982}
Queille, J.-P., Sifakis, J.:
Specification and verification of concurrent systems in CESAR.
In: \textit{Proceedings of the 5th International Symposium on Programming}, Lecture Notes in Computer Science, Bd.~137, S.~337--351. Springer, Berlin, 1982.

\bibitem{holzmann2003}
Holzmann, G.\,J.:
\textit{The SPIN Model Checker: Primer and Reference Manual}.
Addison-Wesley, Boston, 2003.

\bibitem{burch1992}
Burch, J.\,R., Clarke, E.\,M., McMillan, K.\,L., Dill, D.\,L., Hwang, L.\,J.:
Symbolic model checking: $10^{20}$ states and beyond.
\textit{Information and Computation}, 98(2):142--170, 1992.

\bibitem{cimatti2002}
Cimatti, A., Clarke, E.\,M., Giunchiglia, E., Giunchiglia, F., Pistore, M., Roveri, M., Sebastiani, R., Tacchella, A.:
NuSMV~2: An OpenSource Tool for Symbolic Model Checking.
In: \textit{Proceedings of the 14th International Conference on Computer Aided Verification (CAV)}, Lecture Notes in Computer Science, Bd.~2404, S.~359--364. Springer, Berlin, 2002.

\bibitem{mcdermott1998}
McDermott, D., Ghallab, M., Howe, A., Knoblock, C., Ram, A., Veloso, M., Weld, D., Wilkins, D.:
PDDL -- The Planning Domain Definition Language, Version 1.2.
Technical Report CVC TR-98-003/DCS TR-1165, Yale Center for Computational Vision and Control, 1998.

\bibitem{fox2003}
Fox, M., Long, D.:
PDDL2.1: An Extension to PDDL for Expressing Temporal Planning Domains.
\textit{Journal of Artificial Intelligence Research}, 20:61--124, 2003.

\bibitem{helmert2006}
Helmert, M.:
The Fast Downward Planning System.
\textit{Journal of Artificial Intelligence Research}, 26:191--246, 2006.

\bibitem{blum1997}
Blum, A.\,L., Furst, M.\,L.:
Fast Planning through Planning Graph Analysis.
\textit{Artificial Intelligence}, 90(1--2):281--300, 1997.

\bibitem{kautz1992}
Kautz, H.\,A., Selman, B.:
Planning as Satisfiability.
In: \textit{Proceedings of the 10th European Conference on Artificial Intelligence (ECAI)}, S.~359--363, 1992.

\bibitem{hoffmann2001}
Hoffmann, J., Nebel, B.:
The FF Planning System: Fast Plan Generation Through Heuristic Search.
\textit{Journal of Artificial Intelligence Research}, 14:253--302, 2001.

\bibitem{fikes1971}
Fikes, R.\,E., Nilsson, N.\,J.:
STRIPS: A New Approach to the Application of Theorem Proving to Problem Solving.
\textit{Artificial Intelligence}, 2(3--4):189--208, 1971.

\bibitem{vf2}
Cordella, L.\,P., Foggia, P., Sansone, C., Vento, M.:
A (Sub)Graph Isomorphism Algorithm for Matching Large Graphs.
\textit{IEEE Transactions on Pattern Analysis and Machine Intelligence}, 26(10):1367--1372, 2004.

\bibitem{ullmann1976}
Ullmann, J.\,R.:
An Algorithm for Subgraph Isomorphism.
\textit{Journal of the ACM}, 23(1):31--42, 1976.

\bibitem{vardi1986}
Vardi, M.\,Y., Wolper, P.:
An Automata-Theoretic Approach to Automatic Program Verification.
In: \textit{Proceedings of the 1st Annual IEEE Symposium on Logic in Computer Science (LICS)}, S.~332--344, 1986.

\bibitem{gerth1995}
Gerth, R., Peled, D., Vardi, M.\,Y., Wolper, P.:
Simple On-the-fly Automatic Verification of Linear Temporal Logic.
In: \textit{Proceedings of the 15th IFIP WG6.1 International Symposium on Protocol Specification, Testing and Verification}, S.~3--18. Chapman \& Hall, London, 1995.

\bibitem{erol1994}
Erol, K., Hendler, J., Nau, D.\,S.:
HTN Planning: Complexity and Expressivity.
In: \textit{Proceedings of the 12th National Conference on Artificial Intelligence (AAAI)}, S.~1123--1128, 1994.


Epp, S.:
\textit{Model Checking mit dem Subgraph Algorithmus: Formale Verifikation in polynomieller Laufzeit}., 2024.

\bibitem{epp2026pddl}
Epp, S.:
\textit{Temporales Model Checking und PDDL Planning: Formale Verifikation von Planungssystemen}., 2024.


\bibitem{henzinger1996}
Henzinger, T.\,A.:
The Theory of Hybrid Automata.
In: \textit{Proceedings of the 11th Annual IEEE Symposium on Logic in Computer Science (LICS)}, S.~278--292, IEEE, 1996.

\bibitem{alur1995}
Alur, R., Dill, D.\,L.:
A Theory of Timed Automata.
\textit{Theoretical Computer Science}, 126(2):183--235, 1995.

\bibitem{hairer1993}
Hairer, E., N{\o}rsett, S.\,P., Wanner, G.:
\textit{Solving Ordinary Differential Equations~I: Nonstiff Problems}, 2.~Aufl.
Springer, Berlin, 1993.

\bibitem{girard2007}
Girard, A.:
Approximate Bisimulation Relations for Switched Linear Systems.
In: \textit{Proceedings of the 2007 American Control Conference}, S.~4900--4905, IEEE, 2007.

\bibitem{thrun2005}
Thrun, S., Burgard, W., Fox, D.:
\textit{Probabilistic Robotics}.
MIT Press, Cambridge, MA, 2005.

\bibitem{siciliano2009}
Siciliano, B., Sciavicco, L., Villani, L., Oriolo, G.:
\textit{Robotics: Modelling, Planning and Control}.
Springer, London, 2009.

\end{thebibliography}

\end{document}

